% DER_VIII-X: Grothendieck's text with the commentary set in, then the overview, inside a group that applies the chunk's own definitions
\dgrchapter{VIII}{Chapter VIII}{\texorpdfstring{1-types d'homotopie relatifs : leur intégration\,\dots}{1-types d'homotopie relatifs : leur intégration}}
\dgrpart{0}{VIII-X}{Manuscript pages 1--9}{VIII-X\enspace\textperiodcentered\enspace Chapter VIII\enspace\textperiodcentered\enspace manuscript pages 1--9}{Chapter VIII \textperiodcentered\ \textup{VIII-X}}
\dgrchunkbegin
% DER_VIII-X: the definitions of the chunk's preamble, applied inside its group
\let\printedpage\relax\newcommand{\printedpage}[1]{\marginpar{\footnotesize p.~#1}}
\let\msp\relax\newcommand{\msp}[1]{\marginpar{\footnotesize [#1]}}
\let\Grp\relax\newcommand{\Grp}{\underline{\mathrm{Grp}}}
\let\uHom\relax\newcommand{\uHom}{\underline{\mathrm{Hom}}}
\let\uGrp\relax\newcommand{\uGrp}{\underline{\mathrm{Grp}}}
\let\uuGrp\relax\newcommand{\uuGrp}{\underline{\underline{\mathrm{Grp}}}}
\let\fl\relax\newcommand{\fl}{\longrightarrow}
\let\Ch\relax\newcommand{\Ch}{\mathrm{Ch}}
\let\Hot\relax\newcommand{\Hot}{\mathrm{Hot}}
\let\Hom\relax\newcommand{\Hom}{\mathrm{Hom}}
\let\HotO\relax\newcommand{\HotO}{\mathrm{Hot}_0^{\wedge}}
\let\MOP\relax\newcommand{\MOP}{\mathcal{M}_{0\,\Pi}}
\let\M\relax\newcommand{\M}{\mathcal{M}}
\let\uCh\relax\newcommand{\uCh}{\underline{\mathrm{Ch}}}
\let\Homext\relax\newcommand{\Homext}{\mathrm{Hom\,ext}}
\let\Lotimes\relax\newcommand{\Lotimes}{\overset{\mathbf{L}}{\otimes}}
\let\cercle\relax\newcommand{\cercle}[1]{\tikz[baseline=(c.base)]\node[draw,circle,inner sep=1pt](c){#1};}
\let\Chu\relax\newcommand{\Chu}{\underline{\mathrm{Ch}}}
\let\D\relax\newcommand{\D}{\mathbf{D}}
\let\Cat\relax\newcommand{\Cat}{\mathrm{Cat}}
\let\Fl\relax\newcommand{\Fl}{\operatorname{Fl}}
\let\Ob\relax\newcommand{\Ob}{\operatorname{Ob}}
\let\bDelta\relax\newcommand{\bDelta}{\boldsymbol{\Delta}}

% ---- from Der-ChVIII-567_p001-005.tex
% ===================== page 1 =====================
\printedpage{}
% head: (aucun ; page de titre non numérotée)

\begin{center}
{\Large\textsf{LES DÉRIVATEURS}}

\bigskip

\textbf{ALEXANDRE GROTHENDIECK}

\bigskip\bigskip

\textbf{Chapitre VIII}

\medskip

\textbf{1-types d'homotopie relatifs : leur intégration \dots}
\end{center}

\clearpage

% ===================== page 2 =====================
\printedpage{}
% head: (aucun ; page non numérotée)

Ce texte a été déchiffré et transcrit en \LaTeXe{} par M. Künzer. Il a été édité par M. Künzer, J. Malgoire, et G. Maltsiniotis. La transcription est aussi fidèle que possible au manuscrit. Pour les quelques corrections évidentes, ou rares commentaires des éditeurs, ainsi que pour la numérotation originale des pages du manuscrit, les caractères de machine à écrire \texttt{[typewriter]} entre crochets sont utilisés. Un point d'interrogation entre crochets signifie que l'on n'est pas sûr du mot qui précède.

Cette édition est provisoire. Les remarques, commentaires et corrections sont bienvenus. Envoyer un message à :

\textbf{maltsin@math.jussieu.fr}

\hfill G. Maltsiniotis\hspace*{2cm}

\clearpage

% ===================== page 3 =====================
\printedpage{1}
% head: A. Grothendieck : Dérivateurs, Ch. VIII, Édition des Universités de Montpellier II et Paris VII
% foot: Version 2 février 2002

\msp{page 1}

\section*{1\quad Intégration des $1$-types d'homotopie relatifs sur une catégorie.}

Les types d'homotopie relatifs sur une (petite) catégorie $S$ sont définis par les \textit{champs} (en catégories) sur $S^{\wedge}$, qui correspondent essentiellement aux contrafoncteurs $S^o \longrightarrow (\mathrm{Cat})$, ou encore aux catégories $X$ sur $S$ qui sont Cat-fibrées sur $S$. (Dans $(\mathrm{Cat}/S)_{\mathrm{fib}}$, on inverse les flèches qui sont des Hot-équivalences sur les fibres.) D'autre part, les 1-types d'homotopie relatifs sur $S$ doivent se décrire par des \textit{champs en groupoïdes} sur $S$. Si on a un type d'homotopie relatif, défini par
\[ X \longrightarrow S, \qquad\qquad \text{catégorie Cat-fibrée sur } S, \]
ou par
\[
\begin{aligned}
S^o \quad &\longrightarrow \quad \mathrm{Cat}\\
s \quad &\longmapsto \quad X_s\ ,
\end{aligned}
\]
le 1-type d'homotopie associé est défini par
\[ s \longmapsto \Pi(X_s), \qquad\qquad S^o \longrightarrow (\mathrm{Grp}), \qquad\qquad \text{catégorie des groupoïdes,} \]
où $\Pi(X_s)$ est le groupoïde fondamental de $X_s$. On note la catégorie fibrée correspondante sur $S$ par \fbox{$\Pi(X/S)$}. Il est vrai que si
\[ X \ \longrightarrow\ X' \]
est un homomorphisme de catégories fibrées sur $S$ qui induise une Hot-équivalence sur les fibres, elle induit une équivalence fibrée
\[ \Pi(X/S) \ \longrightarrow\ \Pi(X'/S), \]
mais non un isomorphisme. Il y aurait donc du travail à faire pour définir un foncteur
\[
\underbrace{\mathrm{Hot}(S)}_{\text{\footnotesize\begin{tabular}{c}catégorie des\\ types d'homotopie\\ relatifs sur $S$\end{tabular}}}
\quad\xrightarrow{\ \xi\,\longmapsto\,\Pi(\xi/S)\ }\quad
\underbrace{\Grp(S)}_{\text{\footnotesize\begin{tabular}{c}catégorie des\\ groupoïdes\\ fibrés sur $S$\\ (ou catégories\\ des champs en\\ groupoïdes\\ sur $S^{\wedge}$)\end{tabular}}}
\quad .
\]

Mais ce n'est pas pour l'instant mon propos. Plutôt, de donner une formule de transitivité en homologie non commutative, analogue à $R(gf)_* \simeq Rg_*Rf_*$ en cohomologie.

\dgrblock{dgrb418}{Why the target is a 2-category}
\begin{dcomment}{exposition}{p1:Mais ce n'est pas pour l'instant mon propos.}{Why the target is a 2-category}
Grothendieck wants integration to retain the automorphisms of a homotopy type, and not merely its connected components. A groupoid is exactly the amount of structure needed for a homotopy $1$-type. Its objects represent points, its arrows represent homotopy classes of paths, and the nerve has no homotopy groups above degree one. For a category $C$, the groupoid $\Pi C$ is its localization at \emph{all} arrows. In particular, it is generally very different from the maximal subgroupoid of $C$.

The relative version must remember transport as well. We write $\mathsf{Gpd}(S)$ for the $2$-category of categories fibred in groupoids over $S$, functors over $S$, and natural transformations over $S$. Every arrow of a category fibred in groupoids is cartesian; consequently every functor over the base between such categories preserves cartesian arrows, and every $2$-arrow is invertible. A choice of cleavage gives a pseudofunctor $S^{\mathrm{op}}\to\mathsf{Gpd}$, including the isomorphisms comparing successive inverse images. This is the language of \citet[Expos\'e VI, \S\S6--8]{SGA1}. No topology imposing additional descent has been specified here: ``champs'' is being used through the equivalence with fibred categories on the presheaf topos.

Taking fundamental groupoids fibre by fibre is compatible with these comparison isomorphisms. It therefore gives a relative truncation, rather than a mere collection of unrelated groupoids. A weak equivalence of fibres becomes an equivalence of groupoids, which explains Grothendieck's hesitation about a functor to an ordinary category of actual groupoids: localization requires its inverted arrows to become \emph{isomorphisms}. One can either pass to the homotopy category of $\mathsf{Gpd}(S)$, whose morphisms are isomorphism classes of cartesian functors, or retain the entire $2$-category. In modern language the underlying operation is the pointwise truncation
\[
 \operatorname{Fun}(NS^{\mathrm{op}},\mathcal S)
 \longrightarrow
 \operatorname{Fun}(NS^{\mathrm{op}},\mathcal S_{\leq1}).
\]
Here $\mathcal S$ denotes spaces and $\mathcal S_{\leq1}$ their full subcategory of $1$-types. The composition law for integration will be an equivalence with coherent comparison isomorphisms, precisely the form in which it can be useful again under a further change of base.

The terminology continues the programme of \citet[\S69, especially printed pp.~3--5]{Grothendieck1983}: integration is the left adjoint to change of the indexing category. The present calculation isolates a particularly concrete part of that programme, where ordinary categories and their groupoid completions suffice to construct the adjoint.
\end{dcomment}


\clearpage

% ===================== page 4 =====================
\printedpage{2}
% head: A. Grothendieck : Dérivateurs, Ch. VIII, Édition des Universités de Montpellier II et Paris VII
% foot: Version 2 février 2002

\msp{page 2}

Si
\[ f : S \longrightarrow T \]
est un foncteur, on va définir un foncteur
\[ f_! :\ \Grp(S) \ \longrightarrow\ \Grp(T) \]
(image directe de 1-types d'homotopie relatifs). Mais notons d'abord que si
\[ p :\ X \ \longrightarrow\ S \]
est une $S$-catégorie \textit{quelconque} (pas nécessairement Cat-fibrée), on lui associe néanmoins un type d'homotopie relatif, via la catégorie fibrée $\tilde{X}^{(S)}$, ou $\tilde{X}$, ayant comme fibres les $s\backslash X$ (en notant que $s' \longrightarrow s$ dans $S$ définit bien $s\backslash X \longrightarrow s'\backslash X$) (\footnote{Ayant $\tilde{X}$, on définit $\Pi(X/S) \overset{\text{déf}}{=} \Pi(\tilde{X}/S)$.}). On a un $S$-foncteur canonique
\[ X \xrightarrow{\ i\ } \tilde{X}, \qquad\qquad x \longmapsto (x, p(x), \underbrace{p(x)}_{=\,\text{‘}\,s\,\text{’}} \xrightarrow{\ \mathrm{id}\ } p(x)). \]
C'est une $S$-Hot-équivalence si $X$ était déjà fibrée. En fait, sur chaque fibre c'est un homotopisme car on définit alors $\tilde{X}_{s_0} \xrightarrow{\ j\ } X_{s_0}$ de façon évidente, de façon que
\[
X_{s_0} \underset{j}{\overset{i}{\rightleftarrows}} \tilde{X}_{s_0}, \qquad ji \simeq \mathrm{id}, \qquad ij \underbrace{\xrightarrow{\text{homotopie simple}}}_{\text{\footnotesize\begin{tabular}{l}NB cette flèche\\ n'est pas \texttt{[un]}\\ isomorphisme\\ en général.\end{tabular}}} \mathrm{id}.
\]
Revenons alors à $f : S \longrightarrow T$, et considérons un groupoïde fibré sur $S$,
\[ G \ \longrightarrow\ S, \]
ce n'est pas une catégorie fibrée sur $T$, mais on peut lui associer
\[ t \longmapsto t\backslash G, \qquad\qquad T^o \longrightarrow (\mathrm{Cat}), \]
d'où une catégorie fibrée sur $T$, dont on prend le $\Pi$ relatif. Mais il est plus simple de prendre directement
\[ t \ \longmapsto\ \Pi(t\backslash G), \]
qui définit une catégorie fibrée en groupoïdes sur $T$.

\dgrblock{dgrb419}{The comma category computes integration}
\begin{dcomment}{exposition}{p2:qui définit une catégorie fibrée en groupoïdes sur}{The comma category computes integration}
The comma construction makes the variance visible. For $p:X\to S$, an object of $s\backslash_S X$ is a pair $(x,\alpha:s\to p(x))$; an arrow to $(x',\alpha')$ is an arrow $v:x\to x'$ with $p(v)\alpha=\alpha'$. Precomposing $\alpha$ gives the contravariant dependence on $s$. Thus the displayed formula is meaningful even when $p$ has no cartesian lifts.

When $p$ is fibred, those lifts explain the comparison with the actual fibre. Choose a normalized cleavage and send $(x,\alpha)$ to $\alpha^*x\in X_s$. If $v$ is a comma arrow, its composite with the chosen cartesian arrow $\alpha^*x\to x$ factors uniquely through $\alpha'^*x'\to x'$. This defines a functor $j_s:s\backslash_S X\to X_s$. Its composite with the inclusion $i_s$ is the identity, and the cartesian arrows give $i_sj_s\Rightarrow\mathrm{id}$. Applying the nerve yields a homotopy equivalence. The last natural transformation need not be an isomorphism, exactly as the manuscript warns; after applying $\Pi$, it becomes one. This also proves the asserted equivalence $H\simeq\Pi(H/S)$ when $H$ is fibred in groupoids.

For $f:S\to T$ and $G\to S$, the same construction gives
\[
 (f_!G)_t=\Pi(t\backslash_T G).
\]
The subscript on the comma is essential: its objects use arrows $t\to f(s)$ in $T$. Equivalently, if $F:S^{\mathrm{op}}\to\mathsf{Gpd}$ describes $G$, then
\[
 (f_!F)(t)\simeq
 \tau_{\leq1}\operatorname*{colim}_{(s,\alpha)\in(t\backslash f)^{\mathrm{op}}}F(s),
\]
where the colimit on the right is the homotopy colimit in spaces. The total category is the concrete way to form this colimit before truncation. In particular, the $s\backslash G$ printed in the unit on manuscript page~4 must be read as $f(s)\backslash_T G$; this is a notational slip in the edition, not an error introduced by the transcription.

For orientation, let $S=[0\to1]$, and let $F$ be the diagram $F(1)\to F(0)$. Integration over $S$ gives $F(0)$ up to equivalence, whereas its groupoid of sections is $F(1)$. Indeed, the construction above sends an object over $1$ to its inverse image over $0$, and its cartesian arrow becomes invertible in $\Pi G$. Sections are instead determined by their value over $1$. The two endpoints already separate homology from cohomology.
\end{dcomment}


\clearpage

% ===================== page 5 =====================
\printedpage{3}
% head: A. Grothendieck : Dérivateurs, Ch. VIII, Édition des Universités de Montpellier II et Paris VII
% foot: Version 2 février 2002

\msp{page 3}

Ainsi on trouve bien un foncteur (en fait, même un 2-foncteur de 2-catégories)
\[ f_! :\ \Grp(S) \ \longrightarrow\ \Grp(T). \]
Si
\[ u :\ G \ \longrightarrow\ G' \]
est un homomorphisme de groupoïdes fibrés sur $S$, qui induise des Cat-équivalences sur les fibres, alors
\[ f_!(u) :\ f_!(G) \ \longrightarrow\ f_!(G') \]
est aussi une Cat-équivalence sur les fibres. (Cela est formel, à partir du fait que $f_!$ est un 2-foncteur, compte tenu que si on a une équivalence sur les fibres, alors il y a un quasi-inverse -- c'est une histoire générale de catégories fibrées \dots)

Soit maintenant
\[ p :\ X \ \longrightarrow\ S, \]
et prenons le cas $G = \Pi(X/S)$. Je dis qu'on a alors une $T$-équivalence
\[ \boxed{\ f_!(\Pi(X/S)) \ \simeq\ \Pi(X/T),\ } \tag*{$(*)$} \]
en d'autres termes, pour connaître $\Pi(X/T)$ (à $T$-équivalence près), il suffit de connaître $\Pi(X/S)$, et de lui appliquer le foncteur $f_!$ d'intégration des 1-types d'homotopie relatifs, pour $f : S \longrightarrow T$.

Notons que l'on a trivialement
\[ \Pi(X/S) \ =\ p_!(e_X), \]
où $e_X$ désigne le $X$-groupoïde fibré trivial (toutes les fibres ponctuelles). Donc la formule précédente s'écrit comme une formule de transitivité
\[ f_!(p_!(\Gamma)) \ \simeq\ (fp)_!(\Gamma), \]
\msp{page 4}
où $\Gamma$ est le groupoïde fibré trivial sur $X$. Cela suggère de généraliser la formule $(*)$ en
\[ \boxed{\ (fp)_!(\Gamma) \ \simeq\ f_!(p_!(\Gamma))\ } \]
pour \textit{tout} groupoïde fibré sur $X$, dans une situation
\[ X \xrightarrow{\ p\ } S \xrightarrow{\ f\ } T. \]

Ce sera formel, si on établit une propriété de 2-adjonction pour le foncteur $f_!$ :

% ---- from Der-ChVIII-567_p006-010.tex
\printedpage{4}
% head: A. Grothendieck : Dérivateurs, Ch. VIII, Édition des Universités de Montpellier II et Paris VII
% foot: Version 2 février 2002

\textbf{Théorème.} \textit{Soit $f : S \fl T$ une flèche dans} (Cat)\textit{, $G$ un champ en groupoïdes sur $S$, $H$ un champ en groupoïdes sur $T$. On a alors une équivalence de catégories canonique}
\[
\boxed{\ \uHom_{\uGrp(S)}(G, f^*H) \ \simeq\ \uHom_{\uGrp(T)}(f_!G, H).\ }
\]

Les flèches d'adjonction
\[
\left\{
\begin{array}{lclcrcl}
f_!f^*H & \fl & H & \qquad\qquad & H & \in & \mathrm{Ob}\,\uGrp(T)\\[4pt]
f^*f_!G & \longleftarrow & G & \qquad\qquad & G & \in & \mathrm{Ob}\,\uGrp(S)
\end{array}
\right.
\]
devraient se définir sans problème. Pour la deuxième, il faut définir pour $s \in S$
\[
G_s \ \fl\ (f^*(f_!(G)))_s = (f_!(G))_{f(s)} \overset{\text{déf}}{=} \Pi(s\backslash G),
\]
or on a
\[
G_s \fl s\backslash G,
\]
d'où
\[
\underbrace{\Pi(G_s)}_{\begin{subarray}{l} =\ G_s\\ \text{car } G_s\\ \text{est un}\\ \text{groupoïde}\end{subarray}} \ \fl \Pi(s\backslash G),
\]
O.K. C'est fonctoriel en $s$ \dots

\msp{page 5}
Pour la flèche $f_!f^*H \fl H$ sur $T$, il faut trouver pour tout $t \in \mathrm{Ob}\,T$, et fonctoriellement en $t$,
\[
\underbrace{(f_!f^*H)_t}_{\overset{\text{déf}}{=}\Pi(t\backslash f^*(H))} \quad \fl \quad \underbrace{H_t}_{=\,\Pi(H_t)} \ ,
\]
il suffirait donc de définir
\[
t\backslash f^*(H) \ \fl\ H_t,
\]
mais on a le diagramme cartésien
\[
\begin{tikzcd}[column sep=large, row sep=large]
f^*H = H \times_S T \arrow[r] \arrow[d] & H \arrow[d] \\
S \arrow[r, "f"] & T \ni t
\end{tikzcd}
\]
\texttt{[plutôt} $H \times_T S$\texttt{]} induisant
\[
f^*H \ \fl\ H,
\]
\printedpage{5}
% head: A. Grothendieck : Dérivateurs, Ch. VIII, Édition des Universités de Montpellier II et Paris VII
% foot: Version 2 février 2002
d'où
\[
t\backslash f^*H \ \fl\ t\backslash H \qquad\quad \text{(au lieu de s'envoyer dans $H_t$),}
\]
mais d'autre part, comme $H$ est fibrée sur $T$, donc lisse sur $T$, l'inclusion
\[
H_t \fl t\backslash H
\]
est une Hot-équivalence. Ou disons plutôt : $f^*H \fl H$ est un morphisme de catégories sur $T$, induisant
\[
\underbrace{\Pi(f^*H/T)}_{\overset{\text{déf}}{=}\ f_!(f^*H)} \ \fl\ \Pi(H/T),
\]

\msp{page 6}
or le $T$-foncteur canonique
\[
H \fl \Pi(H/T) \qquad\quad (= \text{catégorie sur } T \text{ définie par } t \longmapsto \Pi(t\backslash H))
\]
est une équivalence, parce que $H$ est un fibré en groupoïdes sur $T$. Donc on peut trouver un $T$-quasi-inverse
\[
\Pi(H/T) \fl H
\]
et en composant on trouve
\[
f_!(f^*H) \ \fl\ \Pi(H/T) \ \fl\ H
\]
la flèche cherchée.

Ayant les 2-flèches d'adjonction, je m'arrête là dans les vérifications !

\dgrblock{dgrb420}{Completing the adjunction and its transitivity}
\begin{dcomment}{exposition}{p6:Ayant les 2-flèches d'adjonction}{Completing the adjunction and its transitivity}
The proposed unit and counit are the right ones. To finish their verification without making incompatible choices of quasi-inverses, it is useful to prove a slightly more general universal property. For any category $X\to T$, put $R_TX=\Pi(X/T)$. If $H\to T$ is fibred in groupoids, restriction along $X\to R_TX$ induces an equivalence
\[
 \operatorname{Hom}_T(R_TX,H)\simeq\operatorname{Fun}_T(X,H).
 \tag{VIII.1}
\]
Both sides are groupoids of functors and transformations over $T$.

Here is a direct construction of the inverse. Given $A:X\to H$ over $T$, choose cartesian lifts in $H$ and send $(x,\alpha:t\to p(x))$ to $\alpha^*A(x)$. For a comma arrow $v$, the cartesian factorization of $A(v)$ gives an arrow between these two objects in $H_t$. All such arrows are invertible, so the functor extends uniquely across the localization $t\backslash_T X\to\Pi(t\backslash_T X)$. The comparison isomorphisms of the cleavage supply its dependence on $t$; their cocycle equation follows from the uniqueness of cartesian factorization. A transformation $A\Rightarrow B$ extends by applying $\alpha^*$ to its component at $x$. Thus this construction is a functor on the Hom-groupoids, not just a bijection between classes of objects.

Restricting the extension to $(x,\mathrm{id}_{p(x)})$ recovers $A$. Conversely, an over-$T$ functor $B:R_TX\to H$ sends the cartesian arrow from $(x,\alpha)$ to $(x,\mathrm{id}_{p(x)})$ to a cartesian arrow in $H$. Cartesian factorization therefore gives a natural isomorphism between $B(x,\alpha)$ and $\alpha^*B(x,\mathrm{id})$. This identifies $B$ with the extension of its restriction, naturally also in transformations. It proves (VIII.1), including all the coherence needed for an adjunction.

Apply (VIII.1) to the total category of $G$ viewed over $T$. An over-$T$ functor $G\to H$ is exactly an over-$S$ functor $G\to f^*H$. We obtain the claimed equivalence of Hom-groupoids. The unit sends $g\in G_s$ to the class of $(g,\mathrm{id}_{f(s)})$. The counit sends the class of $((s,h),\alpha:t\to f(s))$ to $\alpha^*h$. The first triangle evaluates the latter formula at an identity arrow. The second uses the canonical cartesian arrow just described, so is the natural isomorphism between extension after restriction and the original functor. Their compatibility is the uniqueness assertion in the same cartesian factorization. Choosing different cleavages changes these formulas by a unique compatible comparison; it does not change the biadjunction.

Transitivity now follows at the same level of precision. For $X\xrightarrow{p}S\xrightarrow{f}T$,
\[
\begin{aligned}
 \operatorname{Hom}_T(f_!p_!\Gamma,H)
 &\simeq\operatorname{Hom}_S(p_!\Gamma,f^*H)\\
 &\simeq\operatorname{Hom}_X(\Gamma,p^*f^*H)\\
 &\simeq\operatorname{Hom}_T((fp)_!\Gamma,H).
\end{aligned}
\]
The bicategorical Yoneda argument identifies the two representing objects. For three composable maps, both parenthesizations represent the same composite restriction; the representing equivalences consequently satisfy the associativity constraint. Taking $\Gamma=e_X$ gives $p_!e_X=\Pi(X/S)$ and the first boxed transitivity formula. What the source leaves unchecked is therefore a coherence verification, not a missing hypothesis in the theorem.
\end{dcomment}


\printedpage{6}
% head: A. Grothendieck : Dérivateurs, Ch. VIII, Édition des Universités de Montpellier II et Paris VII
% foot: Version 2 février 2002

\section*{2\quad Questions de changement de base}

\[
\begin{tikzcd}[column sep=large, row sep=large]
S & S' \arrow[l, "q"'] \arrow[d, "f'"] \\
T \arrow[from=1-1, "f"'] & T' \arrow[l, "p"]
\arrow[from=2-2, to=1-1, bend right=50, shorten <=18pt, shorten >=18pt, "\gamma"']
\end{tikzcd}
\qquad\qquad\qquad
p^*f_!(G) \longleftarrow f'_!q^*(G)
\tag*{$(*)$}
\]

La flèche est définie, pourvu qu'on ait commutativité ou une flèche de commutativité $\gamma : pf' \fl fq$ (pas nécessairement \texttt{[un]} isomorphisme). Par adjonction, il revient au même de donner
\[
\begin{tikzcd}[row sep=large]
p_!f'_!q^*G \arrow[r] \arrow[d, dash, "\wr"', "\text{transitivité}"] & f_!(G) \\
(pf')_!q^*G \arrow[d, "\gamma"] & \\
(fq)_!q^*G \arrow[d, dash, "\text{transitivité}\ \wr"'] & \\
f_!(q_!q^*)G. \arrow[uuur, dashed, "\text{car } q_!q^*G \longrightarrow G"'] &
\end{tikzcd}
\]

On se demande quand la flèche est \texttt{[un]} isomorphisme.

\msp{page 7}
Le théorème suivant doit être vrai, et résulte de mes fourbis généraux.

% vertical bar printed in the left margin beside the following theorem
\textbf{Théorème.} \textit{Si le carré $(*)$ est} cartésien \textit{dans} (Cat)\textit{, $\gamma$ un isomorphisme, et si on suppose que $f$ soit} 1\textit{-lisse, ou $p$} 1\textit{-propre, alors le $T'$-foncteur de changement de base est une équivalence de $T'$-catégories.}

Les ‘fourbis’ généraux seraient en termes d'un dérivateur $\mathbf{D}$ ad hoc, celui des 1-types d'homotopie relatifs. J'y reviendrai par la suite.

\dgrblock{dgrb421}{Base change for truncated homotopy types}
\begin{dcomment}{later}{p7:Les ‘fourbis’ généraux seraient}{Base change for truncated homotopy types}
Grothendieck's next step asks whether integration can be performed after restricting the base. The natural transformation $\gamma:pf'\Rightarrow fq$ fixes the direction of the comparison: contravariance gives $q^*f^*\Rightarrow f'^*p^*$, whose mate is $f'_!q^*\Rightarrow p^*f_!$. At $t'$, its effect before groupoid completion is the evident functor sending a lift $\delta:t'\to f'(s')$ to the composite $p(t')\xrightarrow{p\delta}pf'(s')\xrightarrow{\gamma_{s'}}fq(s')$.

The theorem is valid for the derivator of $1$-types. More generally, set
\[
 \mathbb D_n(I)=h\operatorname{Fun}(NI^{\mathrm{op}},\mathcal S_{\leq n}),
 \qquad n\geq-1.
\]
Its left Kan extensions are homotopy left Kan extensions followed by $n$-truncation; its right Kan extensions are homotopy limits. The truncation is an accessible reflector, and the truncated category remains presentable \citep[Proposition~5.5.6.18 and Remark~5.5.6.21]{Lurie2009HTT}, so these Kan extensions exist for every small indexing category. An $n$-aspherical category is one whose $n$-truncation is a point, with the case $n=-1$ meaning nonemptiness. Smoothness and properness must be taken relative to this theory. They are conditions on the relevant comma categories and their changes of base, not merely on the homotopy types of the strict fibres.

A precise later formulation is \citet[Theorem~1.2, Corollary~4.16, Example~4.17, and \S5]{Maltsiniotis2005Asphericite}. Theorem~1.2 identifies smoothness with the asphericity of the fibre-to-comma inclusions; \S5 treats the corresponding base-change morphism. With his covariant convention replaced by the contravariant convention here, its two alternatives are exactly that the integrating map $f$ is smooth or that the restricting map $p$ is proper. Applying the statement to the localizer detected by fundamental groupoids gives equivalences of groupoids at every $t'$, hence an equivalence of fibred groupoids. No information above the first homotopy group is needed for this particular comparison.

This distinction between a correct integration theorem and an insufficient local-constancy criterion will matter on manuscript page~9. Truncating the \emph{values} of a diagram does not truncate all the coherence by which the values are assembled over its base.
\end{dcomment}


On se pose la question quand le 2-foncteur associé à $f : S \fl T$
\[
f^* :\ \uuGrp(T) \ \fl\ \uuGrp(S)
\]
\printedpage{7}
% head: A. Grothendieck : Dérivateurs, Ch. VIII, Édition des Universités de Montpellier II et Paris VII
% foot: Version 2 février 2002
est 0-fidèle, 1-fidèle, 2-fidèle (\footnote[2]{La $i$-fidélité de $f^*$ ($i \in \{0,1,2\}$) équivaut à celle des foncteurs (obtenus en prenant $G = e_T$)
\[
\underline{\Gamma}(H/T) \ \fl\ \underline{\Gamma}(f^*H/S).
\]}). Si $G$, $H$ sont des fibrés en groupoïdes sur $T$, on considère
\[
\uHom_{\uuGrp(T)}(G,H) \ \fl\ \uHom_{\uuGrp(S)}(f^*G, f^*H),
\]
et suivant que ce foncteur est toujours fidèle, ou pleinement fidèle, ou une équivalence de catégories, on dit que $f^*$ est 0-, 1-, ou 2-fidèle. Dans le cas 2-fidèle, cela implique que si $K$ est un fibré en groupoïdes sur $S$, s'il est équivalent à un $f^*G$, alors $G$ est déterminé à $T$-équivalence près, elle même déterminée à isomorphisme unique près.

\dgrblock{dgrb422}{Sections alone do not test full faithfulness}
\begin{dcomment}{examples}{p7:elle même déterminée à isomorphisme unique près.}{Sections alone do not test full faithfulness}
The three faithfulness conditions distinguish the lifting of transformations, of functors, and of their isomorphisms. The footnote proposes testing them using section groupoids alone. As an assertion for an arbitrary $f$, that reduction is false.

Let $T=[0\to1]$, $S=\{1\}$, and let $f$ be the inclusion. For every contravariant groupoid diagram $H=(H_1\to H_0)$, evaluation at $1$ is an equivalence $\Gamma(T,H)\simeq H_1$: a section consists of $h_1$, $h_0$, and an isomorphism $h_0\simeq u^*h_1$, and the choices of the latter pair form a contractible groupoid. Thus
\[
 \Gamma(T,H)\longrightarrow\Gamma(S,f^*H)
\]
is an equivalence for every $H$.

Nevertheless, take $G_1=H_1=\varnothing$, $G_0=*$, and $H_0=B(\mathbf Z/2)$. The only structure maps from the empty groupoid exist uniquely. A functor $G\to H$ is now just a functor $*\to B(\mathbf Z/2)$, and its natural automorphism group is $\mathbf Z/2$. Consequently
\[
 \operatorname{Hom}_T(G,H)\simeq B(\mathbf Z/2)
 \longrightarrow
 \operatorname{Hom}_S(f^*G,f^*H)\simeq *
\]
is not faithful. The proposed section test therefore fails even for $0$-faithfulness, despite satisfying the strongest of its own three tests.

There is a useful exact replacement. If $\pi_G:G\to T$ denotes the \emph{total category} of the first coefficient, then
\[
 \operatorname{Hom}_T(G,H)
 \simeq\Gamma(G,\pi_G^*H).
\]
Indeed, a section of $G\times_T H\to G$ is precisely a functor $G\to H$ over $T$, and the same identification holds for transformations. After restriction, the base becomes $S\times_T G$. Testing sections on these changed bases, with the pulled-back coefficients, is equivalent to testing the required Hom-groupoids. The missing base change cannot in general be suppressed by invoking an internal Hom: arbitrary restriction of presheaves does not preserve that internal Hom.
\end{dcomment}


\dgrblock{dgrb423}{Factorizations detect all three levels}
\begin{dcomment}{development}{p7:elle même déterminée à isomorphisme unique près.}{Factorizations detect all three levels}
Grothendieck's faithfulness question can be answered without imposing smoothness or properness. The appropriate test keeps both endpoints of an arrow of $T$, as did the factorization categories discussed in the commentary on Chapter~II, manuscript page~53. In the untruncated setting, \citet[arXiv version, Definition~6.9 and Theorem~6.10]{GrothStovicek2016Trees} give the corresponding criterion for homotopical epimorphisms. We spell out its elementary three-level refinement for groupoids. We did not find this combined formulation in the sources checked; the argument uses standard density and adjunction, and carries no claim of priority.

For an arrow $\beta:t\to t'$ of $T$, let $\mathcal F_f(\beta)$ have objects
\[
 (s,a:t\to f(s),b:f(s)\to t'),\qquad ba=\beta,
\]
and arrows $v:s\to s'$ satisfying $f(v)a=a'$ and $b'f(v)=b$.

\textbf{Proposition (the three tests).}\ours\quad
For $i\in\{0,1,2\}$, restriction $f^*:\mathsf{Gpd}(T)\to\mathsf{Gpd}(S)$ is $i$-faithful if and only if every $\mathcal F_f(\beta)$ is $(i-1)$-connected. Explicitly, these conditions are nonemptiness for $i=0$, nonempty connectedness for $i=1$, and $\Pi\mathcal F_f(\beta)\simeq*$ for $i=2$.

\emph{Proof.} The counit makes the factorization categories appear. Write $h_{t'}=\operatorname{Hom}_T(-,t')$, viewed as a discrete groupoid-valued presheaf. The comma formula gives, at $t$,
\[
 (f_!f^*h_{t'})_t
 \simeq\coprod_{\beta:t\to t'}\Pi\mathcal F_f(\beta)
 \longrightarrow \operatorname{Hom}_T(t,t').
 \tag{VIII.3}
\]
The counit maps every summand to its indexing arrow. The biadjunction identifies restriction on $\operatorname{Hom}_T(G,H)$ with precomposition by $\varepsilon_G:f_!f^*G\to G$.

To prove necessity at a fixed $t$, use a coefficient $H=t_*K$, where $t_*$ is the right biadjoint to evaluation at $t$ and $K$ is an arbitrary groupoid. The adjunction reduces the test for $G=h_{t'}$ to precomposition along (VIII.3), with values in $K$. If a summand is empty, take $K=B(\mathbf Z/2)$: an automorphism at that index is invisible after restriction, so faithfulness fails. If a summand has two connected components, the same $K$ allows different automorphisms on those components, which cannot descend from the single index; full faithfulness fails. If a summand is connected but has nontrivial fundamental group, take $K$ to be that summand's groupoid. Its identity functor, extended by constant functors on the other summands, is not isomorphic to a functor factoring through the discrete target. Essential surjectivity fails. These observations give precisely the three necessary conditions.

For sufficiency, first keep $G=h_{t'}$ and allow arbitrary $H$. Under the nonemptiness condition, a modification between two maps $h_{t'}\to H$ is determined by its pullback along $\varepsilon_{h_{t'}}$: at every $t$ and every $\beta$, there is an object over $\beta$ on which to evaluate it. Under connectedness, a modification between the pulled-back maps has equal components on all objects over a fixed $\beta$, since the two maps send every arrow of that groupoid to an identity. It therefore defines a unique component at $\beta$. These components are natural in $t$: reindex an object lying over $\beta$, and use naturality of the given modification. Thus the modification descends uniquely, proving full faithfulness. In the third case, (VIII.3) is itself an equivalence at every $t$, so precomposition is an equivalence of Hom-groupoids.

To pass from representables to an arbitrary $G$, one needs their coherent density, not just generation of connected components. Index the representables by the total category of $G$, sending $(t',g)$ to $h_{t'}$. Their homotopy colimit in groupoid-valued presheaves, evaluated at $t$, is $\Pi(t\backslash_T G)$, hence is equivalent to $G_t$ by the cartesian-lift argument already proved. Thus $G$ is this homotopy colimit. Restriction preserves it, because such colimits are computed pointwise. Taking maps into $H$, respectively into $f^*H$, changes these colimits into homotopy limits of the corresponding Hom-groupoids. Faithful functors remain faithful under these limits, because a modification is determined by its components. Fully faithful functors remain fully faithful, because the components lift uniquely and their compatibility follows by faithfulness. Equivalences remain equivalences by the universal property of a homotopy limit. Equivalently, the homotopy fibres of faithful and fully faithful functors of groupoids are, respectively, $0$-truncated and $(-1)$-truncated spaces; these conditions are preserved by limits. The representable tests therefore imply the required assertion for every $G,H$. This completes the proof.

At level zero the condition has a particularly concrete form: every $t$ is a retract of some $f(s)$. Indeed, apply nonemptiness to $\mathrm{id}_t$; conversely such a retraction supplies a factorization of every arrow with target $t$. For $f:Z\to e$, the sole factorization category is $Z$. The three tests become, respectively, that $Z$ is nonempty, connected, or $1$-connected. Even the strongest condition concerns maps between constant coefficients; it leaves room for the nonconstant diagram over $S^2$ constructed below.
\end{dcomment}


\msp{page 8}
Soit $i \in \mathbf{N}$. On dit que $f$ est $i$-asphérique si $\forall t \in T$, $Z = S/t$ est $i$-asphérique, i.e.
\[
\pi_0(Z),\ \dots,\ \pi_i(Z) \qquad\quad \text{triviaux.}
\tag*{$(*)$}
\]

On définit aussi $(-1)$-asphérique : les $S/t$ sont non-vides. De même pour la notion de $i$-coasphérique.

\dgrblock{dgrb424}{The valid descent assertions in Theorem 3}
\begin{dcomment}{exposition}{p8:De même pour la notion}{The valid descent assertions in Theorem 3}
Theorem~3 contains a useful distinction: full faithfulness of restriction describes how a coefficient already on $T$ is recovered, whereas its essential image describes which coefficients on $S$ descend. The hypotheses in parts~1 and~2 do give these conclusions. The assertion in part~3 needs an additional degree of connectivity.

We give the argument for the valid assertions. Put $n=i-1$. Under hypothesis~a, the inclusions $S_t\to S/t$ are $n$-coaspherical; under hypothesis~b, the inclusions $S_t\to t\backslash S$ are $n$-aspherical. This is the fibre-to-comma criterion for properness or smoothness, respectively. Combining it with the stated asphericity or coasphericity shows that every strict fibre $S_t$ is $n$-connected. The same criterion says that the categories of lifts of an arrow of $T$, with its source fixed in the proper case or its target fixed in the smooth case, are $n$-connected.

For $i=0$, every fibre is nonempty. A natural transformation between two functors $G\to H$ over $T$ is determined by its components over each object of $T$; choosing one object of $S_t$ detects each component. Restriction is therefore faithful on every Hom-groupoid.

For $i=1$, every fibre is connected. Suppose a transformation between the restrictions of two functors $G\to H$ has been given. At $g\in G_t$, its component calculated using $s\in S_t$ is independent of $s$, because naturality along the vertical arrows of $S_t$, and then along zigzags of such arrows, identifies the components. These components are natural within $G_t$. To verify naturality across an arrow $a:t\to t'$, choose a lift of $a$ with one endpoint fixed; its existence follows from the lift-category criterion. Naturality upstairs is the required naturality downstairs. The descended transformation is unique because every fibre is nonempty. This proves full faithfulness on Hom-groupoids.

For $i=2$, base change computes the relevant adjunction on each strict fibre. Under hypothesis~b, if $H$ is on $T$, the counit has the fibrewise form
\[
 (f_!f^*H)_t\simeq\Pi(S_t)\times H_t
 \longrightarrow H_t.
\]
Under hypothesis~a, use the right biadjoint $f_*$ instead: its unit is
\[
 H_t\longrightarrow(f_*f^*H)_t
 \simeq\operatorname{Fun}(\Pi(S_t),H_t).
\]
The right biadjoint is the groupoid of cartesian sections over the comma category $S/t$; the proper base-change formula replaces that comma category by $S_t$. Since $\Pi(S_t)\simeq*$, the displayed arrows are equivalences. The adjunction equivalences on Hom-groupoids now prove that $f^*$ is $2$-faithful. Notice that this argument uses constant coefficients \emph{on each fibre}; it does not assert that every locally constant coefficient there is constant.

The essential-image statement follows by applying the other adjunction arrow. Under hypothesis~b, its restriction to a fibre is the unit
\[
 G|_{S_t}\longrightarrow
 \operatorname{const}_{S_t}\!\left(\Pi\!\int_{S_t}G\right).
\]
If $G|_{S_t}$ is equivalent to a constant groupoid $K$, integration gives $\Pi(S_t)\times K\simeq K$, so this unit is an equivalence. Thus $G\to f^*f_!G$ is fibrewise an equivalence, and $G$ descends. Under hypothesis~a, use $f^*f_*G\to G$ and the same argument with sections: $\operatorname{Fun}(\Pi(S_t),K)\simeq K$. Necessity is immediate, because a coefficient pulled back from $T$ restricts to a constant coefficient on each strict fibre. This proves part~2 with precisely its printed fibrewise-\emph{constant} condition. Its replacement by equivalence of transport functors requires a separate argument and is where the next difficulty lies.
\end{dcomment}


En fait, pour $i \in [-1,+\infty[ \ \subseteq \mathbf{Z}$, on définit un dérivateur $\mathbf{D}_i$ des $i$-types d'homotopie relatifs (cf. plus bas). D'où les notions d'objet $\mathbf{D}_i$-asphérique ($=$ $i$-asphérique), de $\mathbf{D}_i$-équivalence, de foncteur $\mathbf{D}_i$-asphérique et $\mathbf{D}_i$-coasphérique, de foncteur $\mathbf{D}_i$-lisse et $\mathbf{D}_i$-propre, de $\mathbf{D}_i$-fibration et de $\mathbf{D}_i$-fibration \textit{triviale} (i.e. à fibres $\mathbf{D}_i$-asphériques). On doit avoir le

\textbf{Théorème 3.} \textit{Soit $f : S \fl T$ dans} (Cat)\textit{, $i \in \{0,1,2\}$. Supposons l'une des conditions suivantes satisfaite :}

\begin{itemize}[leftmargin=3em]
\item[\textit{a)}] \textit{$f$ est $(i-1)$-propre, et $(i-1)$-asphérique (s'exprime par la condition que certaines catégories, telles les $S/t$, soient $(i-1)$-asphériques).}
\item[\textit{b)}] \textit{$f$ est $(i-1)$-lisse, et $(i-1)$-coasphérique.}
\end{itemize}

\textit{Alors :}

\begin{itemize}[leftmargin=3em]
\item[\textit{1\textsuperscript{o})}] \textit{Le foncteur $f^*$ image inverse}
\[
f^* : \uGrp(T) \ \fl\ \uGrp(S)
\]
\textit{est $i$-fidèle.}
\item[\textit{2\textsuperscript{o})}] \textit{Dans le cas $i = 2$ ($f$ pleinement fidèle), on}
\printedpage{8}
% head: A. Grothendieck : Dérivateurs, Ch. VIII, Édition des Universités de Montpellier II et Paris VII
% foot: Version 2 février 2002
\msp{page 9}
\textit{a de plus ceci : si $G \in \mathrm{Ob}\,\uGrp(S)$ est un fibré en groupoïdes sur $S$, alors $G$ est équivalent à un fibré en groupoïdes de la forme $f^*(H)$ si et seulement si pour tout $t \in T$, la restriction $G_t$ de $G$ à $S_t$ est équivalente à un fibré en groupoïdes constant sur $S_t$. (}\textbf{NB} \textit{La condition a) ou b) implique déjà que $S_t$ est} 1\textit{-connexe, i.e. que son groupoïde fondamental est équivalent à $e$.)}
\item[\textit{3\textsuperscript{o})}] \textit{Soit $Z$ un objet} 1\textit{-connexe dans} (Cat)\textit{, $G$ un fibré en groupoïdes sur $Z$. Pour que $G$ soit $Z$-équivalent à un groupoïde constant sur $Z$, il faut et il suffit que pour toute flèche $u$ dans $Z$, le morphisme de groupoïdes $u^*$ induit sur les fibres de $G$ soit une $Z$-équivalence} (\footnote[3]{On dit alors que $G$ est \textit{essentiellement localement constant} sur $Z$.}).
\end{itemize}

\textbf{Corollaire.} \textit{La condition nécessaire et suffisante de 2\textsuperscript{o}) s'écrit aussi ainsi : pour toute flèche $u$ de $S$ ‘verticale’ sur $T$, i.e. telle que $f(u)$ soit une flèche identique, le foncteur $u^*$ induit par $u$ sur les fibres de $G$ soit une équivalence de groupoïdes.}

\dgrblock{dgrb425}{A finite obstruction in degree two}
\begin{dcomment}{examples}{p9:le foncteur $u^*$ induit par}{A finite obstruction in degree two}
Grothendieck wants a practical descent test that can be checked on the arrows of the fibres. The obstruction to this last simplification already appears over a finite poset, with every fibre the same two-element automorphism group. It makes the distinction from the preceding valid argument quite explicit. The general coherence issue was isolated in the commentary on Chapter~VI, manuscript pages~7 and~74; here it produces a counterexample to a stated theorem.

Let $Z$ be the poset of nonempty proper subsets of $\{0,1,2,3\}$, ordered by inclusion. It has fourteen objects. Its nerve is the barycentric subdivision of the boundary of a tetrahedron, so $|NZ|\cong S^2$ and $\Pi Z\simeq*$. Write $A=\mathbf F_2$. On the strictly increasing triples in $Z$, define
\[
 c(\sigma,\tau,\upsilon)=
 \begin{cases}
 1,&(\sigma,\tau,\upsilon)=(\{0\},\{0,1\},\{0,1,2\}),\\
 0,&\text{otherwise}.
 \end{cases}
\]
Extend $c$ by zero to triples with a repeated object. There is no strictly increasing chain of four objects, so this normalized $2$-cochain is a cocycle.

It is useful to verify directly that it is not a coboundary. Let $z$ be the sum, with coefficients in $\mathbf F_2$, of the twenty-four flags vertex $\subset$ edge $\subset$ face of the tetrahedral boundary. Every edge of this order complex occurs in exactly two such flags: a vertex--edge pair extends to two faces, a vertex--face pair has two intervening edges, and an edge--face pair has two possible vertices. Hence $\partial z=0$, whereas $\langle c,z\rangle=1$. If $c=\delta b$, this pairing would equal $\langle b,\partial z\rangle=0$, a contradiction.

Now define a category $E_c$ over $Z$. There is one object over each $\sigma$, and for $\sigma\subseteq\tau$ the arrows over that inclusion are labelled by $A$; there are no other arrows. The identity has label zero, and composition is
\[
 (\tau\subseteq\upsilon,b)\circ(\sigma\subseteq\tau,a)
 =\bigl(\sigma\subseteq\upsilon,
        a+b+c(\sigma,\tau,\upsilon)\bigr).
 \tag{VIII.2}
\]
The cocycle equation is exactly associativity; normalization gives the unit equations. Each fibre is $BA$. Every arrow is cartesian: composition with a fixed arrow carries the set of lifts of any preceding inclusion bijectively onto the set of lifts of its composite, since in (VIII.2) it adds a fixed element of $A$. Thus $E_c\to Z$ is a category fibred in groupoids. Choosing the lift labelled zero for each inclusion makes every inverse-image functor $BA\to BA$ the identity. Only the comparison isomorphisms for successive inverse images are nontrivial, and they are encoded by $c$. Passing to an equivalent split fibration gives the same counterexample in terms of a strict groupoid-valued functor; the obstruction does not arise from a choice of presentation.

Suppose that $E_c$ were equivalent over $Z$ to a constant fibred groupoid. Its fibre would have to be equivalent to $BA$, so the constant fibred groupoid would admit a section. An equivalence over $Z$ and its quasi-inverse transport that section to a section of $E_c$. Such a section chooses a label $b(\sigma,\tau)$ on each inclusion; its functoriality says
\[
 b(\sigma,\upsilon)
 =b(\sigma,\tau)+b(\tau,\upsilon)+c(\sigma,\tau,\upsilon).
\]
This is precisely $\delta b=c$, which has just been excluded. Therefore $E_c$ is not constant, although $Z$ is $1$-connected and every transport functor is an equivalence. This disproves part~3.

It also disproves the corollary under its own hypotheses. Take $f:Z\to e$. A functor to $e$ is both a Grothendieck fibration and an opfibration, hence smooth and proper for the localizer of $1$-types. Its sole fibre is the $1$-connected category $Z$, so hypotheses~a and~b for $i=2$ both hold. Every arrow is vertical and satisfies the proposed transport test, but $E_c$ is not in the image of $f^*$. Part~2 itself remains valid: its actual condition demands that $E_c|_Z$ be constant, which it is not.
\end{dcomment}


\dgrblock{dgrb426}{The corrected local-constancy criterion}
\begin{dcomment}{status}{p9:le foncteur $u^*$ induit par}{The corrected local-constancy criterion}
The missing datum is classical $2$-monodromy. A locally constant diagram with fibre a groupoid $K$ determines a classifying map
\[
 |NZ|\longrightarrow B\operatorname{Aut}(K),
\]
where $\operatorname{Aut}(K)$ is the groupoid of self-equivalences of $K$ and natural isomorphisms. The target is a $2$-type: its first homotopy group records equivalence classes of self-equivalences, and its second is $\operatorname{Aut}(\mathrm{id}_K)$. A simply connected base removes the first kind of monodromy but can still support the second. In the finite example it is the nonzero class in $H^2(S^2;\mathbf F_2)$.

A precise classification is given by \citet[arXiv version, Theorem~2.2.5 and \S2.3]{PoleselloWaschkies2005}; their Corollary~2.2.4 gives constancy over a $2$-connected base. Applied to the CW complex $|NZ|$, it proves the corrected statement: if $Z$ is $2$-connected, a groupoid diagram is equivalent to a constant one if and only if all its transport functors are equivalences. Alternatively, over a $1$-connected base one must require its remaining $2$-monodromy to vanish. This is a known obstruction, not a new conjecture or a priority claim. Strengthening the corollary by requiring $2$-connected strict fibres, while retaining the hypotheses of part~2, makes the vertical-arrow criterion sufficient by applying this corrected statement separately to every fibre.
\end{dcomment}
\dgrchapter{IX}{Chapter IX}{\texorpdfstring{Retour sur les catégories de fractions $\mathcal{M}W^{-1}$, comparaison avec Quillen}{Retour sur les catégories de fractions MW-1, comparaison avec Quillen}}
\dgrpart{0}{VIII-X}{Manuscript pages 1--25}{VIII-X\enspace\textperiodcentered\enspace Chapter IX\enspace\textperiodcentered\enspace manuscript pages 1--25}{Chapter IX \textperiodcentered\ \textup{VIII-X}}
% ---- from Der-ChIX-678_p001-005.tex
\printedpage{[non numérotée]}
% head: (aucun titre courant imprimé)

\begin{center}
{\LARGE\textsf{LES DÉRIVATEURS}}

\bigskip

\textbf{ALEXANDRE GROTHENDIECK}

\bigskip\bigskip

\textbf{Chapitre IX}

\textbf{Retour sur les catégories de fractions $\boldsymbol{\mathcal{M}W^{-1}}$,\\
comparaison avec Quillen}
\end{center}

\clearpage

\printedpage{[non numérotée]}
% head: (aucun titre courant imprimé)

Ce texte a été déchiffré et transcrit en \LaTeXe{} par M. Künzer. Il a été édité par M. Künzer, J. Malgoire, et G. Maltsiniotis. La transcription est aussi fidèle que possible au manuscrit. Pour les quelques corrections évidentes, ou rares commentaires des éditeurs, ainsi que pour la numérotation originale des pages du manuscrit, les caractères de machine à écrire \texttt{[typewriter]} entre crochets sont utilisés. Un point d'interrogation entre crochets signifie que l'on n'est pas sûr du mot qui précède.

Cette édition est provisoire. Les remarques, commentaires et corrections sont bienvenus. Envoyer un message à :

\textbf{maltsin@math.jussieu.fr}

\hfill G. Maltsiniotis\hspace*{2cm}

\clearpage

\printedpage{1}
% head: A. Grothendieck : Dérivateurs, Ch. IX, Édition des Universités de Montpellier II et Paris VII
% foot: Version 12 mars 2002

\msp{page 1}

\section*{1\quad Fonctorialité de $\Ch^{\Sigma}_{\infty}(a,b)$ en $a$, $b$ dans $\mathcal{M}\Sigma^{-1}$}

On se donne une catégorie $\mathcal{M}$ avec \texttt{[un]} ensemble admissible de flèches $\Sigma \subseteq \mathrm{Fl}(\mathcal{M})$, d'où (en prenant p.~ex. $\underline{t} = (t_n)$, avec $t_n = (\underbrace{\Delta_+\Delta_-}_{\tau_0})^{2n}$, et $\underbrace{t_{n+1}}_{=\,\tau_0^{2n+2}} \longrightarrow \underbrace{t_n}_{=\,\tau_0^{2n}}$ défini comme $p * \mathrm{id}_{\tau_0^{2n}} * p$, où $p : \tau_0 \longrightarrow 1$ la dégénérescence.) On a alors une théorie des $\underline{\Ch}_{\Sigma,\infty}(\mathcal{M},a,b)$, que je note simplement $\Ch_\infty(a,b)$ en tant qu'objet de Hot. Des accouplements
\[
\Ch_\infty(a,b) \times \Ch_\infty(b,c) \longrightarrow \Ch_\infty(a,c) \qquad \text{dans Hot}
\]
avec (nous allons l'admettre) \textit{associativité} dans Hot. Cela définit donc des accouplements sur les $\pi_0(\Ch_\infty(a,b)) \overset{\text{déf}}{=} \pi_0(a,b)$, qui n'est autre que la composition des flèches dans $\mathcal{M}\Sigma^{-1}$. D'autre part, on a pour $a,b \in \mathrm{Ob}\,\mathcal{M}$
\[
\Hom_{\mathcal{M}}(a,b) \longrightarrow \pi_0(\Ch_\infty(a,b)) = \pi_0(a,b),
\]
compatible avec les compositions. Mais on peut l'ignorer pour l'instant, et regarder $\mathrm{Ob}\,\mathcal{M}$, ou simplement $\mathcal{M}_0$, avec la seule structure
\[
\left\{
\begin{array}{ccccccc}
\mathcal{M}_0 \times \mathcal{M}_0 & \longrightarrow & \Hot & \qquad & (a,b) & \longmapsto & \Ch_\infty(a,b) \\[1ex]
\mathcal{M}_0 \times \mathcal{M}_0 \times \mathcal{M}_0 & \longrightarrow & \mathrm{Fl}\,\Hot & \qquad & (a,b,c) & \longmapsto & (\Ch_\infty(a,b) \times \Ch_\infty(b,c) \underbrace{\longrightarrow}_{\substack{\text{flèche}\\ \text{dans Hot}}} \Ch_\infty(a,c))
\end{array}
\right.
\]
avec condition \texttt{[d']}associativité, et aussi une condition d'unités :
\[
\Ch_\infty(a,a) \text{ est muni d'un } 1_a \in \pi_0\Ch_\infty(a,a) = \pi_0(a,a),
\]
jouant le rôle d'une `identité bilatère', dans le sens suivant. Notons que les accouplements
\[
\begin{array}{ccccc}
\Ch_\infty(a,b) & \times & \Ch_\infty(b,c) & \longrightarrow & \Ch_\infty(a,c) \\[1ex]
P & \times & Q & \longrightarrow & R
\end{array}
\]
\msp{page 2}
définissent
\[
\begin{array}{ccc}
P & \longrightarrow & \underline{\Hom}_{\Hot}(Q,R) \\[1ex]
Q & \longrightarrow & \underline{\Hom}_{\Hot}(P,R),
\end{array}
\]
d'où en passant aux $\pi_0$
\[
\begin{array}{ccc}
\pi_0(P) & \longrightarrow & \Hom_{\Hot}(Q,R) \\[1ex]
\pi_0(Q) & \longrightarrow & \Hom_{\Hot}(P,R),
\end{array}
\]
\printedpage{2}
% head: A. Grothendieck : Dérivateurs, Ch. IX, Édition des Universités de Montpellier II et Paris VII
% foot: Version 12 mars 2002
en d'autres termes,
\[
\begin{array}{ccc}
\pi_0(a,b) & \longrightarrow & \Hom_{\Hot}(\Ch_\infty(b,c),\Ch_\infty(a,c)) \\[1ex]
\pi_0(b,c) & \longrightarrow & \Hom_{\Hot}(\Ch_\infty(a,b),\Ch_\infty(a,c)).
\end{array}
\]

\dgrblock{dgrb427}{Composition before functoriality}
\begin{dcomment}{exposition}{p1:avec (nous allons l'admettre)}{Composition before functoriality}
Grothendieck wants a morphism in the localized category to act on an entire space of paths, not merely on its set of components. This is why he temporarily forgets how the path categories were constructed. Once composition and its units are available, change of either endpoint is a formal operation. The hypothesis being made is substantial: associativity is assumed for the particular spaces $\Ch_\infty(a,b)$. The abstract argument that follows does not establish that hypothesis.

Here is the argument with all its variance visible. Write $C(a,b)$ for the proposed Hom object in a monoidal category $(\mathcal H,\otimes,I)$, and suppose that
\[
 m_{abc}:C(a,b)\otimes C(b,c)\longrightarrow C(a,c),
 \qquad j_a:I\longrightarrow C(a,a)
\]
satisfy associativity and the two unit equations, with the associator and unitors of $\mathcal H$ understood. This is enrichment, with the factors ordered in the direction in which a path is traversed; reversing the tensor convention gives the usual notation of \citet[\S\S1.1--1.3]{Kelly1982}. Symmetry is not required. In particular, the Cartesian signs in the general displays on manuscript page 7 should be read as the specified monoidal product $*$.

The ordinary category $C_0$ has
\[
 C_0(a,b)=\operatorname{Hom}_{\mathcal H}(I,C(a,b)).
\]
For $u:I\to C(a,b)$ and $v:I\to C(b,c)$ its composition is
\[
 I\xrightarrow{\sim}I\otimes I
   \xrightarrow{u\otimes v}C(a,b)\otimes C(b,c)
   \xrightarrow{m_{abc}}C(a,c).
\]
The enriched associativity and unit equations give the ordinary ones after precomposition with $I$. Thus the $j_a$ become the identities of $C_0$. Their uniqueness follows already in the endomorphism monoid: if $j$ and $j'$ are both two-sided units, then $j=j\circ j'=j'$. Existence as units on the Hom objects is a separate assumption; an element that is a unit only after applying $\operatorname{Hom}(I,-)$ need not satisfy the stronger enriched equations.

To change endpoints, let $u:a'\to a$ and $v:b\to b'$ belong to $C_0$. Insert the corresponding maps out of $I$ on the two sides of $C(a,b)$ and compose twice. This gives
\[
 C(u,v):C(a,b)\longrightarrow C(a',b').
\]
For composable $u':a''\to a'$ and $v':b'\to b''$, the two composites $C(u',v')C(u,v)$ and $C(uu',v'v)$ are the two parenthesizations of the same fivefold composition, with $u',u$ on the left and $v,v'$ on the right. Associativity identifies them. The unit equations identify $C(1_a,1_b)$ with the identity. Inserting only the left or only the right element also shows that the two actions commute. We have therefore obtained an ordinary bifunctor
\[
 C_0^{\mathrm{op}}\times C_0\longrightarrow\mathcal H.
\]
Applying $\operatorname{Hom}_{\mathcal H}(I,-)$ gives the ordinary Hom bifunctor, including its action on morphisms, rather than only objectwise bijections.

For $\mathcal H=\Hot$, $I$ is a point and $\operatorname{Hom}(I,-)$ is $\pi_0$. One can obtain the two endpoint actions by currying in the Cartesian closed homotopy category, as Grothendieck does on page 2, but the insertion argument shows that closedness is unnecessary. This economy matters when he passes to ind-objects: it is the monoidal structure and its coherence that the construction uses.
\end{dcomment}


On en déduit en particulier pour $a,b \in \mathcal{M}_0$
\[
\left\{
\begin{array}{cccc}
1) & \pi_0(a,a) & \longrightarrow & \Hom_{\Hot}(\Ch_\infty(a,b),\Ch_\infty(a,b)) \\[1ex]
2) & \pi_0(a,a) & \longrightarrow & \Hom_{\Hot}(\Ch_\infty(b,a),\Ch_\infty(b,a)),
\end{array}
\right.
\tag*{$(*)$}
\]
et un objet $\alpha_a$ de $\pi_0(a,a)$ est unité bilatère si ses deux images sont les homomorphismes identiques dans $\Ch_\infty(a,b)$ resp. $\Ch_\infty(b,a)$. Passant aux $\pi_0$ de ceux-ci, cela implique que $\alpha_a$ est unité bilatère pour la composition des $\pi_0(x,y)$, relativement à $b$. Si on veut que ce soit le cas pour \textit{tout $b$,} cela signifie (prenant $b = a$) que $\alpha_a$ est un objet unité bilatère pour la structure de composition associative (donc de monoïde) dans $\pi_0(a,a)$, donc que $\pi_0(a,a)$ est un monoïde \textit{unitaire,} ayant $\alpha_a$ pour objet unité. L'élément $\alpha_a$ est donc déterminé de façon unique. Et on veut que $\alpha_a$ existe
\msp{page 3}
pour tout $a \in \mathcal{M}_0$, donc $\mathcal{M}_0$, muni des $\pi_0(a,b)$ et de leur composition, devient une \textit{catégorie,} soit $\mathcal{M}_\Pi$; et les flèches identiques de celle-ci doivent jouer le rôle d'unité bilatère également pour les `Hom externes' $\Ch_\infty(a,b)$, à valeurs dans Hot.

On peut dire que pour $a = a_0$ fixé, les $\Ch_\infty(a_0,b)$ ($b$ variable) forment les valeurs d'un \textit{foncteur} covariant
\[
\mathcal{M}_\Pi \longrightarrow \Hot \qquad b \longmapsto \Ch_\infty(a_0,b),
\]
en associant à
\[
b \xrightarrow{\ u\ } b', \qquad u \in \pi_0(b,b'),
\]
la flèche
\[
\Ch_\infty(a,b) \longrightarrow \Ch_\infty(b') % [sic] \Ch_\infty(b') -> \Ch_\infty(a_0,b'), argument manquant (cf. la note des éditeurs qui suit)
\]
\texttt{[plutôt} $\Ch_\infty(a_0,b) \longrightarrow \Ch_\infty(a_0,b')$\texttt{]} correspondante $*_1$ [?]. Que ce soit bien là une loi fonctorielle doit résulter formellement de l'associativité des accouplements sur les $\Ch_\infty$. De même, on a une loi fonctorielle, pour $b$ fixé,
\[
\mathcal{M}^o_\Pi \longrightarrow \Hot, \qquad a \longmapsto \Ch_\infty(a,b),
\]
grâce à l'associativité. En fait, l'associativité impliquera aussi qu'on a un \textit{bifoncteur}
\[
\begin{array}{ccc}
\mathcal{M}^o_\Pi \times \mathcal{M}_\Pi & \longrightarrow & \Hot \\[1ex]
(a\,,\,b) & \longmapsto & \Ch_\infty(a,b),
\end{array}
\]
\printedpage{3}
% head: A. Grothendieck : Dérivateurs, Ch. IX, Édition des Universités de Montpellier II et Paris VII
% foot: Version 12 mars 2002
\msp{page 4}
car si
\[
a' \xrightarrow{\ u\ } a, \qquad b \xrightarrow{\ v\ } b',
\]
on a le diagramme
\[
\begin{tikzcd}[column sep=large, row sep=large]
\Ch_\infty(a,b) \arrow[r, "\substack{\text{comp. avec}\\ v}"] \arrow[d, "\text{comp. avec } u"'] & \Ch_\infty(a,b') \arrow[d, "\text{comp. avec } u"] \\
\Ch_\infty(a',b) \arrow[r, "\substack{\text{comp. avec}\\ v}"] & \Ch_\infty(a',b'),
\end{tikzcd}
\]
qui est commutatif grâce à l'associativité pour
\[
\Ch_\infty(a',a),\ \Ch_\infty(a,b),\ \Ch_\infty(b,b').
\]

\textbf{Corollaire} de cette dépendance bifonctorielle : \textit{Si $u : a' \longrightarrow a$, $v : b \longrightarrow b'$ sont dans $\mathcal{M}_\Pi$, alors la flèche correspondante}
\[
\Ch_\infty(a,b) \longrightarrow \Ch_\infty(a',b') \qquad \textit{dans } \Hot
\]
\textit{est un} isomorphisme.

\dgrblock{dgrb428}{An invertibility hypothesis already restored on page 8}
\begin{dcomment}{examples}{p4:de cette dépendance bifonctorielle}{An invertibility hypothesis already restored on page 8}
The intended application is that weakly equivalent endpoints have the same Hom object. For this application $u$ and $v$ are invertible in the localization. The corollary as printed has omitted this condition; manuscript page 8 states the correct version explicitly.

A discrete example detects the omission without any homotopy theory. Enrich the ordered category $0<1$ in sets, regarded as discrete spaces, by its usual Hom sets. Change the first variable by the noninvertible map $u:0\to1$ and fix the second variable at $0$. The induced map is
\[
 C(1,0)=\varnothing\longrightarrow C(0,0)=\{1_0\}.
\]
This is the endpoint action of $u:0\to1$ with the second variable fixed at $0$, and is not an isomorphism. If $u$ and $v$ are invertible, functoriality supplies the inverse $C(u^{-1},v^{-1})$, proving the repaired corollary.

The following display also omits the opposite on the first factor: its domain is $\mathcal M_\Pi^{\mathrm{op}}\times\mathcal M_\Pi$. Both omissions occur in the edition, so they are retained in the source and entered separately in the audit. They do not invalidate the preceding construction of the bifunctor.
\end{dcomment}


Les mêmes arguments devraient marcher si au lieu des $\Hom\,\mathrm{ext}$ à valeurs dans Hot, on les prenait dans Ind\,Hot, ou dans Cat\,hot, ou dans Ind\,Cat\,hot -- je pense bien sûr aux $\mathrm{CH}_\infty(a,b)$ (dans la théorie de la localisation) dans Ind\,Cat\,hot, qui peuvent être regardés comme provenant d'un bifoncteur (où $\mathcal{M}_\Pi = \mathcal{M}\Sigma^{-1}$)
\[
\mathcal{M}_\Pi \times \mathcal{M}_\Pi \longrightarrow \mathrm{Ind}\,\mathrm{Cat}\,\mathrm{hot}.
\]

\msp{page 5}
Voici l'axiomatisation définitive (?). On remplace Hot par une catégorie

\begin{itemize}[label={$\mathcal{H}$,}]
\item catégorie avec produit $*$ associatif, commutatif (?) (\footnote{On ne doit pas avoir besoin de la commutativité.}), unitaire, et toutes les compatibilités habituelles au besoin,
\end{itemize}

l'objet unité pour $*$ sera noté 1, et on notera
\[
\left\{
\begin{array}{ccc}
\pi_0(\xi) & \overset{\text{déf}}{=} & \Hom_{\mathcal{H}}(1,\xi) \\[1ex]
\pi_0 : \mathcal{H} & \longrightarrow & \mathrm{Ens}.
\end{array}
\right.
\]

Si on a
\[
P * Q \xrightarrow{\ u\ } R,
\]

% ---- from Der-ChIX-678_p006-010.tex
\printedpage{4}
% head: A. Grothendieck : Dérivateurs, Ch. IX, Édition des Universités de Montpellier II et Paris VII
% foot: Version 12 mars 2002

on en déduit
\[
\left\{
\begin{array}{lcl}
\pi_0(P) & \xrightarrow{c_u^{\mathrm{l}}} & \mathrm{Hom}_{\mathcal{H}}(Q,R)\\
\pi_0(Q) & \xrightarrow{c_u^{\mathrm{r}}} & \mathrm{Hom}_{\mathcal{H}}(P,R)
\end{array}
\right.
\]
ainsi :
\[
\begin{array}{rlcl}
\text{Si } \alpha\in\pi_0(P),\ \text{i.e. } 1\xrightarrow{\ \alpha\ }P,\ \text{alors} & \multicolumn{3}{l}{\alpha*\mathrm{id}_Q : Q\longrightarrow P*Q}\\
\text{et} & c^{\mathrm{l}}_u(\alpha) & = & u\circ(\alpha*\mathrm{id}_Q)\\
\text{et dualement} & c^{\mathrm{r}}_u(\beta) & = & u\circ(\mathrm{id}_P*\beta).
\end{array}
\]
Si on a trois objets $P_0$, $P_1$, $P_2$ dans $\mathcal{H}$, des accouplements
\[
\left\{
\begin{array}{rcl}
P_0*P_1 & \xrightarrow{u_1} & Q_0\\
P_1*P_2 & \xrightarrow{u_2} & Q_1\\
P_0*Q_1 & \xrightarrow{u_3} & R\\
Q_0*P_2 & \xrightarrow{u_4} & R
\end{array}
\right.
\]
\[
\begin{tikzcd}
a \arrow[r, "P_0"'] \arrow[rr, bend left=20, "Q_0"] \arrow[rrr, bend left=35, "R"] & b \arrow[r, "P_1"'] \arrow[rr, bend right=20, "Q_1"'] & c \arrow[r, "P_2"'] & d
\end{tikzcd}
\]

\msp{page 6}
avec une relation d'associativité, disant que
\[
\begin{tikzcd}
P_0*P_1*P_2 \arrow[r] \arrow[d] & Q_0*P_2 \arrow[d]\\
P_0*Q_1 \arrow[r] & R
\end{tikzcd}
\]
est commutatif, alors on en déduit deux relations de transitivité et une de commutation, pour les opérations associées aux éléments de $\pi_0(P_0)$, $\pi_0(P_1)$, $\pi_0(P_2)$, qu'on n'explicitera pas, sauf l'une :
\[
\begin{array}{rrrcl}
\multicolumn{5}{l}{\text{Soit } \beta\in P_1,\ \gamma\in P_2,}\\
\text{d'où par } \beta : & c_{u_1}(\beta) : & P_0 & \longrightarrow & Q_0,\\
\text{et par } \gamma : & c_{u_4}(\gamma) : & Q_0 & \longrightarrow & R,\\
\text{et par } \beta *_{u_2}\gamma\in\pi_0(Q_1) : & c_{u_3}(\beta\,\texttt{[}{*_{u_2}}\texttt{]}\gamma) : & P_0 & \longrightarrow & R,
\end{array}
\]

\printedpage{5}
% head: A. Grothendieck : Dérivateurs, Ch. IX, Édition des Universités de Montpellier II et Paris VII
% foot: Version 12 mars 2002
\texttt{[plutôt} $\beta\in\pi_0(P_1)$\texttt{,} $\gamma\in\pi_0(P_2)$\texttt{]}, et bien on a
\[
\boxed{\;c(\beta *_{u_2}\gamma) \;=\; c(\gamma)c(\beta).\;}
\]
Il faut encore définir $\beta *_{u_2}\gamma$, via l'accouplement $P_1*P_2\xrightarrow{u_2}Q_1$, on a en effet
\[
\begin{array}{ccc}
\pi_0(P_1)\times\pi_0(P_2) & \longrightarrow & \pi_0(P_1*P_2)\\
(\beta\,,\,\gamma) & \longmapsto & \beta*\gamma,
\end{array}
\]
et on compose avec
\[
\pi_0(u_2) : P_1*P_2 \longrightarrow Q_1
\]
\texttt{[plutôt} $\pi_0(u_2) : \pi_0(P_1*P_2) \longrightarrow \pi_0(Q_1)$\texttt{]}.

\msp{page 7}
On se donne un ensemble $\mathcal{M}_0$, et
\[
\left\{
\begin{array}{rclccrcl}
\mathcal{M}_0\times\mathcal{M}_0 & \xrightarrow{\mathrm{Ch}} & \mathcal{H} & \text{noté} & & (a,b) & \longmapsto & \mathrm{Ch}(a,b)\ \text{ou}\ \mathrm{Hom\,ext}_{\mathcal{H}}(a,b)\\
\mathcal{M}_0\times\mathcal{M}_0\times\mathcal{M}_0 & \xrightarrow{\text{comp.}} & \mathrm{Fl}\,\mathcal{H} & & & (a,b,c) & \longmapsto & \bigl(\mathrm{Ch}(a,b)\times\mathrm{Ch}(b,c)\longrightarrow\mathrm{Ch}(a,c)\bigr)
\end{array}
\right.
\]
avec deux axiomes

\begin{enumerate}
\item[a)] Associativité pour les compositions, pour $(a,b,c,d)\in\mathcal{M}_0^4$
\[
\begin{tikzcd}
\mathrm{Ch}(a,b)\times\mathrm{Ch}(b,c)\times\mathrm{Ch}(c,d) \arrow[r] \arrow[d] & \mathrm{Ch}(a,c)\times\mathrm{Ch}(c,d) \arrow[d]\\
\mathrm{Ch}(a,b)\times\mathrm{Ch}(b,d) \arrow[r] & \mathrm{Ch}(a,d).
\end{tikzcd}
\]
\item[b)] Existence d'objets unité bilatères
\[
1_a\in\pi_0(\mathrm{Ch}(a,a)) \qquad \text{pour tout } a\in\mathcal{M}_0
\]
\end{enumerate}

Alors les $\pi_0(\mathrm{Ch}(a,b))$ et leur composition définissent sur $\mathcal{M}_0$ comme ensemble d'objets une structure de \textit{catégorie} $\MOP$, ayant les $1_a$ comme identités, et on trouve que
\[
(a,b)\longmapsto\mathrm{Ch}(a,b)
\]
peut s'interpréter à partir d'un \textit{bifoncteur} canonique
\[
\mathcal{M}^o_{0\,\Pi}\times\MOP \longrightarrow \mathcal{H},
\]
dont le composé avec
\[
\mathcal{H}\xrightarrow{\pi_0}\mathrm{Ens}
\]
n'est autre que le bifoncteur $(a,b)\longrightarrow\mathrm{Hom}_{\MOP}(a,b)$,
\printedpage{6}
% head: A. Grothendieck : Dérivateurs, Ch. IX, Édition des Universités de Montpellier II et Paris VII
% foot: Version 12 mars 2002
\msp{page 8}
\[
\mathrm{Hom}_{\MOP}(a,b) \;\simeq\; \pi_0(\mathrm{Ch}(a,b)).
\]
Cette relation est la \textit{définition} de $\underline{\mathrm{Hom}}_{\MOP}$, mais ici la relation est entendue comme une % Reading resolved against PDF leaf 18: a single underline; see corrections.csv.
bijection \textit{fonctorielle} en $a$, $b$.

On trouve notamment que si
\[
u : a'\longrightarrow a, \qquad v : b\longrightarrow b'
\]
sont des flèches \textit{inversibles} dans $\MOP$, alors
\[
\mathrm{Ch}(a,b)\longrightarrow\mathrm{Ch}(a',b')
\]
correspondant dans $\mathcal{H}$ est \texttt{[un]} \textit{isomorphisme.}

On appliquera tout ceci surtout aux deux situations déduites d'une situation $(\mathcal{M},\Sigma)$, avec
\[
\begin{array}{rlcllcl}
1^o) & \mathcal{H} & = & \mathrm{Hot}, & \mathrm{Ch}(a,b) & = & \mathrm{Ch}^{\Sigma}_{\infty}(\mathcal{M};a,b)\\
2^o) & \mathcal{H} & = & \mathrm{Ind\,Cat\,hot}, & \mathrm{Ch}(a,b) & = & \mathrm{CH}^{\Sigma}_{\infty}(\mathcal{M};a,b)\ \text{(\footnotemark[2])}.
\end{array}
\]
\footnotetext[2]{$*$ dans Hot et Ind\,Cat\,hot est le produit cartésien sans plus.}
Dans les deux cas, on a donc un foncteur sur $(\mathcal{M}\Sigma^{-1})^o\times(\mathcal{M}\Sigma^{-1})$, à valeurs soit dans Hot, soit dans Ind\,Cat\,hot. (\textbf{NB} Les deux structures de catégorie sur $\mathcal{M}_0$ déduites des % [sic] des -> de, « déduites des ces deux foncteurs »
ces deux foncteurs sont canoniquement isomorphes.)
\[
\begin{tikzcd}[column sep=small]
 & \mathrm{Hot} & (a,b)\longmapsto\mathrm{Ch}^{\Sigma}_{\infty}(a,b)\\
(\mathcal{M}\Sigma^{-1})^o\times(\mathcal{M}\Sigma^{-1}) \arrow[ur] \arrow[dr] & & \\
 & \mathrm{Ind\,Cat\,hot} & (a,b)\longmapsto\mathrm{CH}^{\Sigma}_{\infty}(a,b)
\end{tikzcd}
\]

\msp{page 9}
Comme les foncteurs canoniques
\[
\begin{tikzcd}
\mathrm{Ind\,Cat\,hot} \arrow[r, "\mathrm{Ind}\,\varphi"] & \mathrm{Ind\,Hot} \arrow[r, "\mathrm{Ind}\,\psi"] & \mathrm{Ind\,Hot}_0^{\wedge} \arrow[r, "\varinjlim"] & \mathrm{Hot}_0^{\wedge},\\
\mathrm{Cat\,hot} \arrow[u, hook] \arrow[r, "\varphi"] & \mathrm{Hot} \arrow[u, hook] \arrow[r, "\psi"] & \mathrm{Hot}_0^{\wedge} \arrow[u, hook] \arrow[ur, "\mathrm{id}"'] &
\end{tikzcd}
\]
commutant aux produits finis, toute théorie de `Hom externes' à valeurs dans une de ces catégories, en définit une à valeurs dans chacune de celles dans laquelle elle s'envoie.
\printedpage{7}
% head: A. Grothendieck : Dérivateurs, Ch. IX, Édition des Universités de Montpellier II et Paris VII
% foot: Version 12 mars 2002
D'ailleurs, ces foncteurs ne sont pas tous pleinement fidèles (c'est en tous cas faux pour $\varphi$ et $\mathrm{Ind}\,\varphi$, et j'ignore ce qui en est de $\psi$ et $\mathrm{Ind}\,\psi$ -- les flèches verticales par contre sont pleinement fidèles (\footnote[3]{et également, $\mathrm{Ind\,Hot}_0^{\wedge}\xrightarrow{\varinjlim}\mathrm{Hot}_0^{\wedge}$ n'est \textit{pas} pleinement fidèle.}).) Mais ils induisent des bijections pour les $\pi_0$. Donc la catégorie $\MOP$ déduite d'une théorie de Hom externes à valeurs dans une de ces catégoires, et celle % [sic] catégoires -> catégories, faute d'orthographe
déduite dans une autre où elle s'envoie, est la même à isomorphisme canonique près. Pour ce qui est de cette définition de $\MOP$, rien n'est changé donc quand on regarde les $\mathrm{CH}_{\infty}$ comme des objets de Ind\,Hot, ou seulement de $\mathrm{Ind\,Hot}_0^{\wedge}$, ou même seulement de $\HotO$, ou quand on regarde les $\mathrm{Ch}_{\infty}$ comme des objets
\msp{page 10}
dans $\HotO$. Mais bien sûr, la donnée des accouplements de composition perd chaque fois en précision (sauf peut-être en passant de Hot à $\HotO$, si $\psi$ est pleinement fidèle \dots\,). Quand on prend les images des $\mathrm{CH}^{\Sigma}_{\infty}(a,b)$ et $\mathrm{Ch}^{\Sigma}_{\infty}(a,b)$ dans $\HotO$ (la dernière catégorie du diagramme), on trouve des objets canoniquement isomorphes (fonctoriellement en $a$, $b$ dans $\MOP\simeq\mathcal{M}\Sigma^{-1}$). Quand on les regarde comme étant respectivement dans Ind\,Hot et dans $\mathrm{Hot}\hookrightarrow\mathrm{Ind\,Hot}$, on trouve un homomorphisme fonctoriel
\[
\mathrm{CH}^{\Sigma}_{\infty}(a,b)\longrightarrow\mathrm{Ch}^{\Sigma}_{\infty}(a,b) \qquad\quad \text{dans Ind\,Hot ou dans } \mathrm{Ind\,Hot}_0^{\wedge} \text{ au choix.} \tag*{$(*)$}
\]
La $\varinjlim$ de $\mathrm{CH}^{\Sigma}_{\infty}(a,b)$ dans Hot n'existe pas en général (sauf si l'ind-objet est essentiellement constant), mais elle existe dans $\HotO$, et l'isomorphisme canonique
\[
\varinjlim_{\HotO}\mathrm{CH}^{\Sigma}_{\infty}(a,b)\xrightarrow{\ \sim\ }\mathrm{Ch}^{\Sigma}_{\infty}(a,b) \qquad\quad \text{dans } \HotO \tag*{$(**)$}
\]
est la flèche déduite de $(*)$ dans $\mathrm{Ind\,Hot}_0^{\wedge}$.

\dgrblock{dgrb429}{What the finite tests forget}
\begin{dcomment}{status}{p9:D'ailleurs, ces foncteurs ne sont pas tous pleinement fidèles}{What the finite tests forget}
Grothendieck is trying to recover a space of morphisms from its values on a manageable collection of diagram shapes. The attraction is clear: if those values determined both spaces and maps between them, a difficult coherence argument could be replaced by calculations on finite diagrams. The notation needs its earlier definition. In Chapter V, manuscript pages 51--52, $\Hot_0$ is the full subcategory of homotopy types represented by categories with finitely many objects and arrows \citep[Chapter V, printed pp.~31--32]{GrothendieckDerivateursV}. It is not the subcategory of finite CW complexes. For example, the one-object category of a finite nontrivial group is finite, whereas its nerve need not have finite dimension.

Choose a small skeleton $\mathcal A$ of $\Hot_0$. The functor in question is
\[
 \psi(X)(K)=[K,X],\qquad K\in\mathcal A.
\]
It preserves products, and evaluation at the one-point category recovers $\pi_0X$. These useful facts survive the negative answer to the full-faithfulness question. The commentary on Chapter~V, manuscript p.~52, discusses the obstruction (``Finite unbased tests miss a symmetric-group shift'' and ``The modern scope of the final obstruction''). The general theorem is that no set of unpointed homotopy types jointly detects isomorphisms in $\Hot$ \citep[Theorem~2.1]{ArlinChristensen2023}. Applied to $\mathcal A$, this says that $\psi$ is not even conservative, and hence is not fully faithful.

There is also a separate answer to the faithfulness question, which is what the hoped-for associativity argument actually needs. The presheaf category on $\mathcal A$ has a faithful set-valued functor
\[
 F\longmapsto\coprod_{K\in\operatorname{Ob}\mathcal A}F(K):
\]
a natural transformation is determined by its components, which this disjoint union retains. If $\psi$ were faithful, the composite would make $\Hot$ concrete. Freyd's nonconcreteness theorem rules this out; his argument already applies to a full subcategory of simply connected, finite-dimensional CW complexes, so forgetting basepoints does not evade the obstruction \citep[pp.~1 and 5 of the 2004 reprint]{Freyd1970}. Thus $\psi$ is not faithful either. These are statements with the usual fixed-universe interpretation of small spaces and small sets; restricting the allowed spaces to a particular small category changes the question.

Consequently, agreement after passage to $\Hot_0^{\wedge}$ cannot in general prove equality of the two composites in $\Hot$. It follows as well, by restriction to constant ind-objects, that $\operatorname{Ind}\psi$ is not fully faithful. None of this invalidates the underlying category of morphisms obtained by taking $\pi_0$; it identifies exactly the additional coherence that manuscript page 13 still has to supply.
\end{dcomment}


\dgrblock{dgrb430}{Three meanings of a limiting space}
\begin{dcomment}{examples}{p10:dans Hot n'existe pas en général}{Three meanings of a limiting space}
Grothendieck keeps an ind-object because it records the stage at which a family of paths can be represented. Passing to its union, or to a colimit of its homotopy types, asks a different question. The phrase ``n'existe pas en général'' should be read as the absence of a general existence theorem here. Essential constancy is a sufficient condition, not a necessary condition for an ordinary colimit in $\Hot$.

For a complete elementary example, let $E_n=\{0,\ldots,n\}$, with the discrete topology and the evident inclusions, and let $E=\mathbf N$. For every space $Y$ there are natural bijections
\[
 [E,Y]=(\pi_0Y)^{\mathbf N}
       \simeq\varprojlim_n(\pi_0Y)^{E_n}
       =\varprojlim_n[E_n,Y].
\]
Hence $E$ is the ordinary colimit in $\Hot$. Nevertheless, the ind-object $(E_n)$ is not essentially constant. If it were isomorphic to a constant object $Z$, its functor of maps to constant objects would imply $Z\simeq E$. An inverse morphism from the constant $E$ to $(E_n)$ factors through some $E_n$, by the defining formula for maps in an ind-category. Its composite with the cocone would make the identity of $\mathbf N$ factor through a finite set on $\pi_0$, which is impossible.

The same example explains the footnote about the colimit functor on presheaves. Let $F_n$ be the constant presheaf with value $E_n$ on $\mathcal A$, and $F$ the constant presheaf with value $\mathbf N$. Then $\operatorname{colim}F_n=F$. The identity of $F$ has no lift to a morphism from the constant ind-object $F$ to $(F_n)$: such a lift would factor through one finite $F_n$, again contradicting evaluation at any object of $\mathcal A$. Thus
\[
 \operatorname{Ind}(\widehat{\mathcal A})
       \longrightarrow\widehat{\mathcal A},\qquad
 (F_n)\longmapsto\operatorname{colim}F_n,
\]
is not full.

A further qualification is needed for the displayed comparison with $\Ch_\infty$. A filtered homotopy colimit $K=\operatorname{hocolim}K_n$ is tested by a \emph{compact} homotopy type $A$ through
\[
 \operatorname{colim}_n[A,K_n]\simeq[A,K].
\]
Finite-category homotopy types are not all compact. Take $A=BC_2=\mathbf{RP}^{\infty}$ and $K_n=\mathbf{RP}^{n}$. Their union is $A$, but the identity of $A$ cannot factor, even up to homotopy, through any $K_n$: it would kill the nonzero mod-$2$ cohomology in degrees above $n$ \citep[Theorem~3.19]{Hatcher2002}. Thus the finite-category definition of $\Hot_0$ alone does not justify the comparison in manuscript page 10. For the particular path system, one must either prove the required factorization on these test categories by its special categorical structure, or impose an appropriate finite-test/compactness hypothesis. We do not replace that missing model-specific argument by a general assertion about filtered homotopy colimits.
\end{dcomment}


\msp{page 11}
Dans les deux cas CH et Ch (qu'on va désigner par le symbole commun $C_{\infty}(a,b)$), la fonctorialité de cet objet $C_{\infty}(a,b)$ (même au sens le plus faible, comme objet de $\HotO$) par rapport à $a$ et $b$ (même au sens le plus faible, comme foncteur $\mathcal{M}^o\times\mathcal{M}\longrightarrow\HotO$ et pas même $(\mathcal{M}\Sigma^{-1})^o\times(\mathcal{M}\Sigma^{-1})\longrightarrow\HotO$) n'est pas évidente du tout sur la définition directe en termes du système inductif explicite dans Cat
\[
C_{\infty}(a,b) \;=\; (\mathrm{Ch}_{\tau_0^{2n}}(a,b))_{n\in\mathbf{N}},
\]
car une flèche $a'\longrightarrow a$ ou $b\longrightarrow b'$ dans $\mathcal{M}$, et même qu'elle serait dans $\Sigma$, ne définit aucunement une flèche dans Cat
\[
\begin{array}{rcl}
C_{\infty}(a,b) & \longrightarrow & C_{\infty}(a',b)\\
\text{ou}\ \ C_{\infty}(a,b) & \longrightarrow & C_{\infty}(a,b').
\end{array}
\]
Cette difficulté rappelle celle qu'avait à surmonter \textsc{Quillen}, dans Homotopical Algebra, pour affronter la différence entre `left homotopies' et `right homotopies'. Je me rappelle comment
\printedpage{8}
% head: A. Grothendieck : Dérivateurs, Ch. IX, Édition des Universités de Montpellier II et Paris VII
% foot: Version 12 mars 2002
\msp{page 12}
on peut néanmoins définir de telles flèches, comme résultant d'une flèche (beaucoup plus riche) de `composition' de chemins entre $a'$ et $a$, et $a$ et $b$, ou entre $a$ et $b$, et $b$ et $b'$, soit en général donc l'accouplement
\[
C_{\infty}(a,b)\times C_{\infty}(b,c)\longrightarrow C_{\infty}(a,c),
\]
dont la définition directe n'est pas moins problématique, et ne peut en tous cas \textit{pas} se donner comme une flèche dans Cat lui-même, ou dans Ind\,Cat. Interprétant p.~ex. $\mathrm{Ch}_{\infty}(a,b)$ en termes de chemins \textit{infinis} (des deux cotés) de $a$ à $b$, en un sens évident (mais % [sic] cotés -> côtés, accent omis
`essentiellement finis' des deux cotés), on peut définir aussi % [sic] cotés -> côtés, accent omis
\[
\mathrm{Ch}^{-}_{\infty}(a,b), \qquad \mathrm{Ch}^{+}_{\infty}(a,b)
\]
comme les sous-catégories pleines formées respectivement de chemins qui sont totalement dégénérés en degrés $\geq 0$ (i.e. $a_0\longleftarrow a_1\longrightarrow a_2\longleftarrow a_3\cdots$ sont tous des identités), ce qui revient aux chemins infinis seulement à gauche, i.e. s'arrêtant à $a_0$, et symétriquement pour $\mathrm{Ch}^{+}_{\infty}(a,b)$. Même tabac pour les $\mathrm{CH}^{+}$, $\mathrm{CH}^{-}$.

\msp{page 13}
Ceci dit, on a ceci (revenant à la notion $C_{\infty}$ ambivalente) :

\begin{enumerate}
\item[a)] les inclusions pleinement fidèles
\[
C^{-}_{\infty}(a,b)\hookrightarrow C_{\infty}(a,b)\hookleftarrow C^{+}_{\infty}(a,b)
\]
sont des isomorphismes dans Ind\,Cat\,hot resp. dans Hot, et
\item[b)] On a bel et bien une composition des chemins
\[
C^{-}_{\infty}(a,b)\times C^{+}_{\infty}(b,c)\longrightarrow C_{\infty}(a,c).
\]
\end{enumerate}

L'\textit{associativité} dans le cas du CH n'offre d'ailleurs pas de difficulté, à cause de la possibilité de se ramener aux
\[
C_m(a,b)\times C_n(b,c)\times C_p(c,d)\rightrightarrows C_N(a,d)
\]
pour $N$ grand, et prouver l'egalité comme flèches dans Cat\,hot -- ce qui résultera du lemme % [sic] egalité -> égalité, accent omis
d'homotopie sur les intervalles dans Cat. Dans le cas du $\mathrm{Ch}_{\infty}$ par contre je n'ai pas trouvé de démonstration. (Faute seulement d'avoir essayé sérieusement de m'y mettre? Je n'ai pas trop envie, vu qu'il se pourrait que $\mathrm{Hot}\longrightarrow\HotO$ soit pleinement fidèle, auquel cas le travail devient inutile \dots\,)

\dgrblock{dgrb431}{The path model and the localization it is meant to represent}
\begin{dcomment}{later}{p13:dans le cas du CH n'offre}{The path model and the localization it is meant to represent}
Grothendieck is trying to make concatenation independent of an arbitrary placement of the middle endpoint. Moving one path to the left and the other to the right produces an actual functor; the homotopy equivalences from these one-sided path categories then give composition in $\Hot$. For the ind-object $\mathrm{CH}_\infty$, every finite collection of paths and reparametrizations can be handled at one larger stage. His interval homotopies therefore give the asserted associativity in $\mathrm{Ind\,Cat\,hot}$. For the realized $\Ch_\infty$, the text explicitly leaves the associativity proof unfinished. An eventual equality at every finite stage is not, by itself, a specified compatible homotopy on the whole realization.

The commentary on Chapter~V, manuscript p.~12 (``What composition in an ind-category actually asserts'') explains the relevant ind-Hom formula. Here it gives the exact associativity argument. Put $\mathcal A=\mathrm{Cat\,hot}$ and write $C_m(a,b)$ for the finite stages. Morphisms from the triple product into $\mathrm{CH}_\infty(a,d)$ are computed by
\[
 \varprojlim_{m,n,p}\ \varinjlim_N
 \operatorname{Hom}_{\mathcal A}\bigl(
 C_m(a,b)\times C_n(b,c)\times C_p(c,d),C_N(a,d)\bigr).
\]
For each fixed $(m,n,p)$, concatenate the three finite words before returning to the chosen standard shape. Parenthesization does not change this concatenated word; it changes only the finite reparametrization and the inserted identity padding. At a common larger stage the interval-homotopy lemma identifies those reparametrizations; its relative endpoint condition is the one explained in ``Even indices and reparametrization relative to the ends'' and ``Sliding an even block gives a natural transformation'' in the commentary on Chapter~V, manuscript pp.~14 and~21. Hence the two representatives have the same class in the displayed filtered colimit for every $(m,n,p)$, which proves equality in the inverse limit. No choice of these finite homotopies compatible as $(m,n,p)$ varies is required for this equality of ind-morphisms. Concatenating an identity path similarly adds only identity padding, so the same argument proves the strong unit equations for $\mathrm{CH}_\infty$. This is the precise advantage of the ind-object in the present argument.

The general target of this construction was already available in a different model. The hammock localization of \citet[\S2; Propositions 2.2 and 3.1]{DwyerKan1980Calculating} is a simplicial category whose composition concatenates and reduces hammocks; its category of components is $\mathcal M[W^{-1}]$. Thus it supplies coherent mapping spaces, with strictly associative composition at that model level, for a relative category. These statements concern the hammock construction. They do not identify its simplicial sets with Grothendieck's fixed-type path categories, whose transformations and stabilization conventions must be compared separately.

This distinction continues the commentary on Chapter~V, manuscript p.~27 (``From components of path categories to mapping spaces''). Equality of $\pi_0$, agreement on finite tests, and a weak equivalence of the mapping spaces are three different assertions. We have not located in the sources checked a comparison theorem stated for exactly the bilateral stabilization and unrestricted natural transformations used here. The model-specific comparison, including compatibility with concatenation, is consequently left identified as an unfinished verification in these pages, not advertised as a new open problem in homotopy theory. The abstract bifunctor construction above is complete once its stated associativity and unit hypotheses hold.
\end{dcomment}


% ---- from Der-ChIX-678_p011-015.tex
\printedpage{9}
% head: A. Grothendieck : Dérivateurs, Ch. IX, Édition des Universités de Montpellier II et Paris VII
% foot: Version 12 mars 2002

\msp{page 14}

\section*{2\quad Comparaison avec Quillen}

Il se pose un certain nombre de questions sur ce foncteur $(\M\Sigma^{-1})^o \times (\M\Sigma^{-1}) \longrightarrow \Hot$, et sa relation aux constructions classiques, notamment à la théorie de \textsc{Quillen}. Je vais maintenant noter
\[
\begin{aligned}
\Ch^{\Sigma}_{\infty}(a,b) \;&=\; \Homext_{\Hot}(a,b) && \text{ou simplement}\\
&\phantom{\;=\;} \Homext(a,b).
\end{aligned}
\]

1$^o$). Soient $f,g : a \rightrightarrows b$ dans $\M$, alors (si $\M$ est une `model category') \textsc{Quilllen} définit un ensemble % [sic] Quilllen -> Quillen, three l's printed
\[
\pi_1(a,b;f,g) \qquad \text{ou} \qquad \pi_1(f,g)
\]
des classes d'homotopie de $f$ à $g$, de telle façon que $\mathrm{Hom}_{\M}(a,b)$ devient l'ensemble des objets d'un groupoïde, dont les flèches sont les éléments de $\pi_1(f,g)$, pour une notion de composition naturelle. Je suspecte que l'on a
\[
\pi_1(f,g) \;=\; \mathrm{Hom}_{\underbrace{\Pi(\uCh^{\Sigma}_{\infty}(a,b))}_{\substack{\text{groupoïde}\\ \text{fodamental}\\ \text{de } \uCh^{\Sigma}_{\infty}(a,b)}}}(f,g), % [sic] fodamental -> fondamental, missing n
\]
quand $f$ et $g$ sont considérés comme des `chemins infinis' (en insérant chaque $f : a \longrightarrow b$ dans
\[
\cdots \underbrace{a_{-2}}_{=a} \xleftarrow{\ \mathrm{id}\ } \underbrace{a_{-1}}_{=a} \xrightarrow{\ f\ } \underbrace{a_{0}}_{=b} \xleftarrow{\ \mathrm{id}\ } \underbrace{a_{1}}_{=b} \longrightarrow \cdots\ \ )
\]

\dgrblock{dgrb432}{Homotopies as arrows in a mapping groupoid}
\begin{dcomment}{exposition}{p14:Je suspecte que l'on a}{Homotopies as arrows in a mapping groupoid}
Grothendieck's question asks for the first layer of the mapping space, where homotopies themselves become morphisms. Its modern formulation begins with a derived mapping space $\operatorname{Map}^{h}(a,b)$, and takes its fundamental groupoid. The objects corresponding to actual maps $a\to b$ form a full subgroupoid. They need not meet every component until the usual cofibrancy and fibrancy hypotheses, or replacements, have been imposed.

For a cofibrant object $a$ and a fibrant object $b$ in a simplicial model category, $\operatorname{Map}(a,b)$ is a Kan complex. A simplicial path from $f$ to $g$ represents a homotopy; two such paths give the same arrow of its fundamental groupoid precisely when they are homotopic relative to their endpoints. Concatenation in the groupoid is defined by filling an inner two-dimensional horn. Different fillers give the same class because the Kan conditions provide the requisite relative homotopy. Constant paths give identities and reversed paths give inverses. Consequently
\[
 \operatorname{Hom}_{\Pi\operatorname{Map}(a,b)}(f,g)
\]
is the exact object Grothendieck wants, not the set of equivalence classes of maps $a\to b$.

Quillen's original construction has precisely the cofibrant-source and fibrant-target qualifications: see Chapter I, \S2, Proposition 1.2.1 in the checked retypesetting; Chapter II, \S2 identifies this groupoid with the fundamental groupoid of the simplicial mapping complex \citep{Quillen1967}. For a general model category, \citet[Propositions 4.4--4.5 and 4.8]{DwyerKan1980Function} compare resolutions with the hammock mapping spaces and with the simplicial enrichment when it exists. These results justify the same description using $\operatorname{Map}^{h}$ after replacements. The identification with $\Pi\Ch^W_\infty(a,b)$ still requires the particular path-model comparison just discussed. The source's conjectural equality should therefore be stated as a natural equivalence of groupoids after that comparison, not as a literal equality of two chosen constructions.
\end{dcomment}


\msp{page 15}

Donc la définition des $\pi_1(f,g)$ (si c'est bien ça) se prolonge de façon naturelle aux chemins infinis (essentiellement finis) arbitraires.

2$^o$) Supposons que $\M$ admette un objet à la fois initial et final $0_{\M}$ ($\M$ `ponctuée'). Je dis qu'alors toutes les catégories $\uCh^{\Sigma}_{\infty}(a,b)$ sont ponctuées de façon naturelle : à savoir, par le chemin infini de $a$ à $b$ défini (comme plus haut) à partir de $f = 0_{a,b} : a \longrightarrow b$, l'application nulle de $a$ à $b$. On a envie alors de regarder les $\uCh^{\Sigma}_{\infty}(a,b)$ comme définissant des éléments dans le $\Hot_{\bullet}$ (`Hot ponctué'). C'est plus fin que de dire qu'on se donne une classe neutre
\[
e \;\longrightarrow\; \Ch(a,b) \qquad \text{dans Hot,}
\]
car le foncteur de $\Hot_{\bullet}$ dans Hot (ou dans $e\backslash\Hot$) n'est pas fidèle. Le problème, c'est de définir les flèches
\[
\Ch(a,b)\times\Ch(b,c) \;\longrightarrow\; \Ch(a,c)
\]
dans $\Hot_{\bullet}$, et pas seulement dans Hot, et avec associativité (et les unités) dans $\Hot_{\bullet}$ également. Si on applique la définition avec $\Ch^{-}(a,b)$ et $\Ch^{+}(b,c)$, le composé des éléments nuls $a \xrightarrow{\ 0\ } b$ de $\Ch^{-}(a,b)$ et $b \xleftarrow{\ \mathrm{id}\ } b \xrightarrow{\ 0\ } c$ de $\Ch^{+}(b,c)$
\printedpage{10}
% head: A. Grothendieck : Dérivateurs, Ch. IX, Édition des Universités de Montpellier II et Paris VII
% foot: Version 12 mars 2002
\msp{page 16}
est l'élément de $\Ch_{\infty}(a,c)$ que voici
\[
\begin{tikzcd}[row sep=huge]
 & a_{-2} & a_{-1} \arrow[l] \arrow[r] & a_{0} \arrow[d, equal] & a_{1} \arrow[l] \arrow[r] & a_{2} & \cdots \\
{}_{\mathrm{id}}\ \cdots & a & a \arrow[l, "\mathrm{id}"'] \arrow[r, "0"'] & b & b \arrow[l, "\mathrm{id}"'] \arrow[r, "0"'] & c & \cdots_{\mathrm{id}} \arrow[l, "\mathrm{id}"']
\end{tikzcd}
\]
ce n'est \textit{pas} l'élément nul de $\Ch_{\infty}(a,c)$ ! Mais on a des flèches canoniques
\[
e^{-}_{a,b}\circ e^{+}_{b,c} \;\;\rightleftarrows\;\; e_{a,c} \qquad \text{dans } \Ch_{\infty}(a,c)
\]
qui sont l'identité sur tous les sommets sauf celui d'indice 0. Il y a un sorite à développer, comme quoi une paire
\[
(f : K_{\bullet} \longrightarrow L_{\bullet}, c)
\]
d'un foncteur entre catégories ponctuées, et d'un chemin $c : f(e_{K_{\bullet}}) \longrightarrow e_{L_{\bullet}}$, définit dans Hot une flèche canonique
\[
\mathrm{hot}_{\bullet}(K_{\bullet}) \;\longrightarrow\; \mathrm{hot}_{\bullet}(L_{\bullet})
\]
avec transitivité si on a
\[
(f' : L_{\bullet} \longrightarrow M_{\bullet}, c' : f'(e_{L_{\bullet}}) \longrightarrow e_{M_{\bullet}}),
\]

\dgrblock{dgrb433}{The path needed to preserve a basepoint}
\begin{dcomment}{exposition}{p16:Il y a un sorite à développer}{The path needed to preserve a basepoint}
Grothendieck wants the zero map to remain a specified point throughout composition, because its loop group will define suspension and loops. A point chosen in the unpointed homotopy category does not contain the required information. For example, based maps $S^1\to BG$ are classified by elements of a group $G$, whereas free homotopy identifies conjugate elements. For $G=S_3$, the two distinct transpositions $(12)$ and $(13)$ give different based classes and the same free class. This proves the asserted failure of faithfulness of the forgetful functor.

The proposed extra datum is exactly the right one at the level of spaces. For pointed CW spaces $(K,k)$ and $(L,l)$, a map $f:K\to L$ together with a path $c$ from $f(k)$ to $l$ defines a morphism in the pointed homotopy category. To see this without suppressing the path, attach an interval to $K$ at $k$ and point the resulting space $K^+$ at its free endpoint. Define $F:K^+\to L$ by $f$ on $K$ and by $c$ on the interval. The collapse $r:K^+\to K$ is a pointed homotopy equivalence. The required morphism is $[F][r]^{-1}$. The same construction applies to nerves of pointed categories after geometric realization, because a vertex inclusion is a cofibration.

For a second pair $(g,d)$, the composite is represented by $(gf,(gc)*d)$. Indeed, the space obtained by attaching two successive intervals to $K$ maps to the target by $f$, then the first path followed under $g$, then $d$. Collapsing the subdivision of the two intervals identifies this representative with the single attached interval carrying $(gc)*d$. Repeating this argument gives associativity: the increasing piecewise-linear reparametrizations of a concatenated interval are homotopic relative to the two endpoints. Homotopic choices of $c$ relative endpoints give the same pointed morphism. Arbitrary paths with the same endpoints need not do so; their difference can act through the fundamental group.

The specific claim immediately before this construction is stronger and is false as written, but it has a simple canonical repair. Denote the displayed concatenation by $P$ and the standard zero path from $a$ to $c$ by $Q$. Then $P$ has $b$ at vertices $0,1$, while $Q$ has $c$ at both. A transformation cannot be the identity at every vertex except $0$: the backward identity edge between $1$ and $0$ forces its components at those two vertices to agree. Instead define $\alpha:P\to Q$ by
\[
 \alpha_0=\alpha_1=0_{b,c},\qquad
 \alpha_i=\mathrm{id}\quad(i\ne0,1).
\]
The square at $-1\to0$ commutes because a composite with a zero map is zero; the square at $1\to0$ compares two identity edges; and the square at $1\to2$ says $\mathrm{id}_c\,0_{b,c}=0_{b,c}$. Every other square is an identity square. Since transformations of paths are unrestricted, this is an admissible canonical comparison for every $W$ containing identities. It supplies the needed basepoint path without a reparametrization or an additional choice. A reverse arrow with zero components would force $\mathrm{id}_c=0_{c,c}$ at the last square and is not available in general. This repairs the single basepoint comparison; associativity and the other coherence conditions for the particular pointed composition still require verification.
\end{dcomment}


\msp{page 17}

d'où
\[
f'' = f'\circ f : K_{\bullet} \longrightarrow M_{\bullet}, \qquad \text{et}\ \ c'' : f''(e_{K_{\bullet}}) \longrightarrow e_{M_{\bullet}}), % [sic] closing parenthesis without opening parenthesis -> (f'' = f'\circ f : ..., c'' : ...), unbalanced parenthesis
\]
définit ainsi
\[
\begin{tikzcd}
f'(f(e_{K_{\bullet}})) \arrow[r, "f'(c)"] \arrow[rr, bend right=12, "{c''\ \text{par définition}}"'] & f'(e_{L_{\bullet}}) \arrow[r, "c'"] & e_{M_{\bullet}}
\end{tikzcd}
\]

\noindent\llap{?\quad}Cela nous donne donc bel et bien des accouplements dans $\Hot_{\bullet}$ (et aussi dans $\mathrm{Cat\,hot}_{\bullet}$?). Il est difficile de croire qu'ils pourraient ne pas être associatifs! D'ailleurs $\Hot_{\bullet}$, avec le produit cartésien ordinaire (\footnote[4]{Non, il faut prendre le produit $\xi_{\bullet} \divideontimes \eta_{\bullet}$ ($\xi_{\bullet}\times\eta_{\bullet}$, avec $\xi_{\bullet}\vee\eta_{\bullet}$ contracté en $e$).}), satisfait à l'axiomatique de tantôt pour $\mathcal{H}$, et $\Hot_{\bullet} \longrightarrow \Hot$ est compatible avec ce produit (\footnote[5]{Non, mais on s'en tire quand même \ldots}), et avec $\pi_0$. % Reading resolved against PDF leaf 22: \divideontimes, meaning smash; see corrections.csv.

\dgrblock{dgrb434}{Why the monoidal unit changes in the pointed theory}
\begin{dcomment}{exposition}{p17:Il est difficile de croire}{Why the monoidal unit changes in the pointed theory}
The two footnotes identify a genuine change of monoidal structure. In pointed spaces the Cartesian unit is the one-point space, and
\[
 \operatorname{Hom}_{\Hot_\bullet}(*,X)=\{*\}.
\]
It cannot recover the components of $X$, so the Cartesian structure does not instantiate the earlier convention $\pi_0(X)=\operatorname{Hom}(I,X)$. With the smash product, the unit is $S^0$, and
\[
 \operatorname{Hom}_{\Hot_\bullet}(S^0,X)=\pi_0(X).
\]
Thus the corrected enrichment uses smash composition
\[
 C(a,b)\wedge C(b,c)\longrightarrow C(a,c).
\]
The identity is the pointed map $S^0\to C(a,a)$ carrying the nonbasepoint to $1_a$ and the basepoint to zero. Composition must kill either zero input, not merely the pair of zero inputs. That is the additional requirement hidden by a single basepoint comparison for a Cartesian product. The standard pointed simplicial model of derived mapping spaces has this absorbing-zero property, hence its composition factors through the smash product. This is the familiar monoidal example in \citet[\S1.1]{Kelly1982}.

Forgetting basepoints is not strong monoidal for smash and Cartesian products. Nevertheless there is a natural quotient map
\[
 U(X)\times U(Y)\longrightarrow U(X\wedge Y)
\]
and a unit map $*\to U(S^0)$ selecting its nonbasepoint. These make $U$ lax monoidal: both associativity composites collapse exactly those tuples with a zero coordinate, and both unit composites leave the other coordinate unchanged. Composing these maps with smash composition recovers associative unpointed composition. The footnote ``on s'en tire quand m\^eme'' can therefore be made precise without asserting preservation of products. The symbol in the edition is the crossed product $\divideontimes$, explained there as collapsing the wedge, and is faithfully transcribed; its meaning here is the smash product.
\end{dcomment}


Si tout ceci marche, on aurait donc une théorie des $\Homext_{\bullet}$ à valeurs dans $\Hot_{\bullet}$; et sans doute aussi pour les $\mathrm{CH}_{\infty}$ comme objets de $\mathrm{Ind\,Hot}_{\bullet}$, et non seulement de $\mathrm{Ind\,Hot}$. Mais
\printedpage{11}
% head: A. Grothendieck : Dérivateurs, Ch. IX, Édition des Universités de Montpellier II et Paris VII
% foot: Version 12 mars 2002
de toutes façons on ne peut en général laisser tomber Cat hot, qui reste importante -- donc il faudrait avoir des objets $\mathrm{CH}_{\infty}(a,b)$ dans $\mathrm{Ind\,Cat\,hot}_{\bullet}$, pour avoir une côte bien taillée.

Les $\pi_0(C(a,b))$ seront à présent des ensembles ponctués, et ce sera par l'image de l'application nulle $a \xrightarrow{\ 0\ } b$, i.e. le composé
\[
a \longrightarrow 0 \longrightarrow b
\]

\msp{page 18}

qu'il faut composer avec l'elément déduit du chemin tronqué % [sic] l'elément -> l'élément, missing accent
\[
a \longrightarrow \underbrace{0}_{=a_0} \longleftarrow b \qquad \ldots
\]

Il faudrait montrer que 0 est aussi un objet neutre de $\M\Sigma^{-1}$, i.e. que
\[
\pi_0(a,0) = \pi_0(0,a) = \{0\} \qquad \text{pour tout } a \in \mathrm{Ob}\,\M
\]
et interpréter l'élément nul dans $\pi_0(a,b)$ comme l'application nulle dans $\M\Sigma^{-1}$. De façon générale, il faut des conditions qui assurent (sans que $\M$ soit forcément ponctuée) que
\[
\M \xrightarrow{\ \gamma\ } \M\Sigma^{-1}
\]
commute aux sommes finies, aux produits finis (donc $\gamma(\varphi_{\M}) = \varphi_{\M\Sigma^{-1}}$, resp. $\gamma(e_{\M}) = e_{\M\Sigma^{-1}}$) (\footnote[6]{à faire.}).

\dgrblock{dgrb435}{Zero objects and finite sums after localization}
\begin{dcomment}{exposition}{p18:Il faudrait montrer que 0 est aussi}{Zero objects and finite sums after localization}
Grothendieck needs the distinguished component to be an actual zero morphism of the localized category. At the level of ordinary localization this requires no additional compatibility condition on $W$. If $z$ is initial in $\mathcal M$, write $i_x:z\to x$ for its unique maps. For every ordinary arrow $u:x\to y$ we have $ui_x=i_y$. If $u$ is inverted, the same equality gives $u^{-1}i_y=i_x$ in the localization. An induction along any finite zigzag starting at $z$ therefore identifies it with $i_x$. Thus $z$ remains initial. The reversed argument preserves terminal objects. In particular a zero object remains zero, and $a\to0\to b$ is the zero morphism in $\mathcal M[W^{-1}]$. This is also the nullary case of descent of adjunctions through localization \citep[Lemma 1.1.6]{Krause2021}.

The displayed path $a\to0\leftarrow b$ needs a different caution: its backward edge belongs to $W$ only when $b\to0$ does. The two forward zero maps $a\to0\to b$ always give the required ordinary zero morphism; inserting backward identity edges puts this composite into the permitted zigzag shapes. The preservation of the zero object therefore does not depend on the inadmissible shortened path or on the unresolved pointed composition.

For binary products there is a useful sufficient condition. Suppose $\mathcal M$ has finite products and $W$ contains identities and is closed under their finite products. The product functor and the diagonal then descend to the localizations, as do the unit and counit of $\Delta\dashv\prod$. Naturality for inverse generators follows by multiplying the original naturality squares by the inverses of their $W$-edges; the triangle identities remain equalities after localization. Finally
\[
 (\mathcal M^n)[(W^n)^{-1}]
 \simeq (\mathcal M[W^{-1}])^n
 \qquad(n<\infty):
\]
a tuple of finite zigzags is represented by one tuple-shaped zigzag after padding the finitely many words by identities, and the relations are coordinatewise. Consequently the descended diagonal still has the descended product as right adjoint. This proves that localization preserves finite products. Reversing arrows gives the corresponding statement for finite coproducts. One can weaken closure by requiring only that the relevant products or coproducts of $W$-arrows become invertible.

For an additive category, closure of $W$ under finite direct sums is already enough for the ordinary localized category to be additive; a five-lemma hypothesis is unnecessary for this conclusion \citep[Lemma 2.2.1]{Krause2021}. Indeed, products and coproducts survive and the original canonical comparison from their sum to their product remains an isomorphism, giving biproducts. Addition is therefore defined by the usual diagonal--biproduct--codiagonal formula. The morphism $-1_x$ remains an inverse for the addition on $x$, and its naturality extends over inverted arrows, so every Hom monoid is a group. A central $k$-action similarly extends. These conclusions concern ordinary Hom groups; constructing the proposed complexes of higher Hom information is the further problem on the next page.
\end{dcomment}


Supposons enfin que $\M$ soit \textit{additive,} et même $k$-linéaire ($k$ un anneau fixé), i.e. on a de plus un homomorphisme d'anneaux
\[
k \;\longrightarrow\; \mathrm{End}(\mathrm{id}_{\M}).
\]

On doit supposer quelque compatibilité de $\Sigma$ à la structure additive, p. ex. que la somme directe de deux flèches dans $\Sigma$ est dans $\Sigma$, et sans doute quelque chose de plus fort, du genre `lemme des cinq' pour $\Sigma$. (Mais on ne suppose pas forcément $\M$ abélienne, donc on aurait dû mal à formuler un lemme des cinq?) % [sic] dû mal -> du mal, accent on du

\msp{page 19}

Mais peut-être ici faut-il un axiome plus subtil, p.ex. (si $\underbrace{\Sigma}_{\text{suspension}}$ et $\Omega$ existent (\footnote[7]{cf. plus bas ce que ça signifie.})) que l'ensemble des morphismes d'adjonction $\Sigma\Omega \longrightarrow \mathrm{id}$, $\mathrm{id} \longrightarrow \Omega\Sigma$ soit \texttt{[formé d']}isomorphismes, ce qui implique que $\Omega$ resp. $\Sigma$ est pleinement fidèle. En tous cas, l'attente naturelle ici, c'est que $\Ch^{\Sigma}_{\infty}(a,b)$ soit non seulement dans $\Hot_{\bullet}$, mais même dans la catégorie $D^{\leq 0}(k\text{-Mod})$, sous-catégorie pleine de la catégorie dérivée $D^{-}(k\text{-Mod})$ formée des $K$ tels que $H^{i}(K) = 0$ si $i > 0$. On peut aussi la noter $D_{\geq 0}(k\text{-Mod})$, l'interprétant en
\printedpage{12}
% head: A. Grothendieck : Dérivateurs, Ch. IX, Édition des Universités de Montpellier II et Paris VII
% foot: Version 12 mars 2002
termes de complexes de `chaînes' plutôt que de cochaînes. Ici les accouplements doivent correspondre à l'opération $* = \Lotimes_k$ dans $D_{\geq 0}(k\text{-Mod})$, admettant $k$ comme objet unité, et
\[
\pi_0(K_{\bullet}) \;\simeq\; H_0(K_{\bullet}).
\]
Donc on doit avoir
\[
\Homext(a,b) \Lotimes_k \Homext(b,c) \;\longrightarrow\; \Homext(a,c),
\]
où cette fois Hom ext désigne l'objet de $D_{\geq 0}(k\text{-Mod})$, au lieu d'un objet de Hot seulement, ou de $\Hot_{\bullet}$. On utilise les foncteurs
\[
D_{\geq 0}(k\text{-Mod}) \longrightarrow \Hot_{\bullet} \longrightarrow \Hot
\]
\msp{page 20}
pour faire le lien avec la théorie ponctuée dans $\Hot_{\bullet}$, ou \texttt{[la]} théorie générale dans Hot. Si on désigne par
\[
\gamma :\ D_{\geq 0}(k\text{-Mod}) \;\longrightarrow\; \Hot_{\bullet}
\]
le premier de ces foncteurs, il ne commute pas aux produits $*$ ($\Lotimes_k$ et $\times$ respectivement), mais on a, pour $P_{\bullet}$, $Q_{\bullet}$ variable dans $D_{\geq 0}(k\text{-Mod})$, une flèche fonctorielle canonique % [sic] variable -> variables, agreement with P, Q
\[
\gamma(P_{\bullet} \Lotimes_k Q_{\bullet}) \xleftarrow{\ \mathit{can.}\ } \gamma(P_{\bullet})\times\gamma(Q_{\bullet})
\]
`$k$-bilinéaire', ce qui montre que de donner des accouplements
\[
P_{\bullet} * Q_{\bullet} \xrightarrow{\ c_k\ } R_{\bullet},
\]
d'où
\[
\gamma(P_{\bullet} * Q_{\bullet}) \xrightarrow{\ \gamma(c_k)\ } \gamma(R_{\bullet}), \tag*{$(*)$}
\]
est l'un et l'autre plus fin que de se donner seulement
\[
\gamma(P_{\bullet}) * \gamma(Q_{\bullet}) \xrightarrow{\ c_{\bullet}\ } \gamma(R_{\bullet}), \tag*{$(*')$}
\]
vu qu'une donnée $(*)$ impliquera une donne $(*')$ par % [sic] donne -> donnée, missing letters
\[
\begin{tikzcd}
 & \gamma(P_{\bullet} * Q_{\bullet}) \arrow[dr, "\gamma(c)"] & \\
\gamma(P_{\bullet}) * \gamma(Q_{\bullet}) \arrow[ur, "\text{can.}"] \arrow[rr, "c_{\bullet}"] & & \gamma(R_{\bullet}).
\end{tikzcd}
\]
Ceci vu, on s'attend bien à avoir une théorie \textit{forte} des accouplements
\printedpage{13}
% head: A. Grothendieck : Dérivateurs, Ch. IX, Édition des Universités de Montpellier II et Paris VII
% foot: Version 12 mars 2002
\msp{page 21}
\[
\Homext_k(a,b) \Lotimes_k \Homext_k(b,c) \;\longrightarrow\; \Homext_k(a,c),
\]
d'où on déduira les accouplements plus faibles dans $\Hot_{\bullet}$ (au lieu de $D_{\geq 0}(k\text{-Mod})$)
\[
\Homext_{\Hot_{\bullet}}(a,b) \times \Homext_{\Hot_{\bullet}}(b,c) \;\longrightarrow\; \Homext_{\Hot_{\bullet}}(a,c).
\]

\dgrblock{dgrb436}{A sufficient differential graded setting}
\begin{dcomment}{later}{p19:l'attente naturelle ici}{A sufficient differential graded setting}
Grothendieck now asks for more than additive Hom groups: he wants composition on their derived complexes before taking $H_0$. A precise sufficient setting is a differential graded model, for example complexes of modules with quasi-isomorphisms. Choose one cofibrant replacement $Qa\to a$ for each object, in the projective model structure; all complexes are fibrant. The complexes $\operatorname{Hom}^\bullet(Qa,Qb)$ then compose literally by differential graded composition, since the middle representative is the same in the two factors. This is the enhancement described in \citet[\S\S2.2 and 3.2]{Keller2006DG}. The claim here is about this specified setting. The text does not state a set of hypotheses on an arbitrary additive relative category that would construct such an enhancement, and the existence of an adjoint equivalence $\Sigma\dashv\Omega$ alone does not give that construction.

For a commutative ring $k$, put
\[
 C(a,b)=\tau_{\leq0}\operatorname{Hom}^\bullet_k(Qa,Qb)
       \simeq\tau_{\leq0}\mathbf R\operatorname{Hom}_k(a,b).
\]
The good truncation retains the negative degrees and the cycles in degree zero. It belongs to $D^{\leq0}(k)$, or equivalently to $D_{\geq0}(k)$ after reversing degrees. Differential graded composition gives
\[
 C(a,b)\mathbin{\overset{\mathbf L}{\otimes}}_k C(b,c)
       \longrightarrow C(a,c),
\]
using the Koszul interchange to put the factors into the usual composition order. Composing two cycles of degree zero gives a cycle, and the sum of two nonpositive degrees is nonpositive, so composition restricts to the good truncations. Replacing the two input complexes by projective resolutions provides the indicated derived tensor map. Its associativity is inherited from differential graded composition and the coherent tensor comparison; its unit is $k\to C(a,a)$ carrying $1$ to the identity. Since $k$ is projective over itself,
\[
 \operatorname{Hom}_{D(k)}(k,C(a,b))=H^0C(a,b)
 =\operatorname{Hom}_{D(k)}(a,b).
\]

The unqualified ``ring'' in the source needs attention at exactly this tensor product. Two left modules over a noncommutative ring do not have the displayed tensor product with values in left modules. One may assume $k$ commutative, as above, or use the appropriate bimodule category. Moreover, the image of $k\to\operatorname{End}(\mathrm{id}_{\mathcal M})$ is central: any two natural endomorphisms commute by applying the naturality of one to a component of the other. Thus the scalar action already factors through a commutative image.

The connective complex produces a pointed space by Dold--Kan and forgetting the simplicial module structure. The convention is
\[
 \pi_n\operatorname{Map}^{h}(a,b)
   =H^{-n}\mathbf R\operatorname{Hom}_k(a,b)
   =\operatorname{Hom}_{D(k)}(a[n],b),\qquad n\geq0,
\]
with the first equality at the zero basepoint for $n>0$. The higher positive Ext groups of two degree-zero modules are not positive homotopy groups of this mapping space. This is the grading distinction explained in the commentary on Chapter~III, manuscript p.~12 (``Keeping all degrees of the mapping complex'').

For a concrete calculation, take $a=k[0]$ and $b=k[1]$, with cohomological convention $K[1]^i=K^{i+1}$. The source is projective, so the ordinary Hom complex already computes derived Hom and is $k[1]$, concentrated in degree $-1$. Dold--Kan gives $K(k,1)$; there is one component, its fundamental group is the additive group of $k$, and every higher homotopy group vanishes. On the other hand $\operatorname{Hom}_{D(k)}(k[1],k[1])=k$, exactly the group predicted by suspension representability. This example makes visible the information that would disappear if one kept only the ordinary zero map $k\to k[1]$.
\end{dcomment}


\dgrblock{dgrb437}{The bilinear comparison is lax monoidal}
\begin{dcomment}{exposition}{p20:une flèche fonctorielle canonique}{The bilinear comparison is lax monoidal}
Grothendieck correctly asks for a comparison map in the direction needed to transport composition. It need not be an isomorphism. Let $P$ and $Q$ be nonnegative chain complexes and choose nonnegative projective resolutions $P'\to P$ and $Q'\to Q$. Their Dold--Kan simplicial modules $A=\Gamma P'$ and $B=\Gamma Q'$ have a degreewise map of pointed sets
\[
 U(A)\times U(B)\longrightarrow U(A\otimes_k B),
 \qquad (a,b)\longmapsto a\otimes b.
\]
The shuffle equivalence identifies the normalized complex of $A\otimes_k B$ with $P'\otimes_k Q'$, which computes $P\otimes_k^{\mathbf L}Q$. This yields the source's comparison in $\Hot_\bullet$. It factors through the smash product because either zero input is sent to zero. For coherence, use the unital associative Alexander--Whitney comparison $\Gamma P'\otimes_k\Gamma Q'\to\Gamma(P'\otimes_k Q')$; its normalized map is inverse to shuffle up to chain homotopy. The derived construction is independent of the chosen projective replacements. Dold--Kan itself is Theorem 14.24.3 of \citet{StacksDoldKan}; the tensor comparisons, their identities, and the induced monoidal structure on $\Gamma$ are checked in \citet[\S2.3, especially equation (2.8)]{SchwedeShipley2003}.

Already with $P=Q=k[0]$, the map on components is
\[
 k\times k\longrightarrow k,\qquad (x,y)\longmapsto xy,
\]
which is not injective for any nonzero ring, since $(0,0)$ and $(0,1)$ have the same image. It nonetheless carries a tensor pairing to a pointed or unpointed pairing by composition. Thus ``more information'' means that the bilinear factorization is specified, not that the forgetful functor is fully faithful or that its monoidal comparison is invertible. The displayed Cartesian product is legitimate for this weaker pairing; to apply the enriched unit convention in pointed spaces one uses its smash factorization, as explained above.
\end{dcomment}


3$^o$) Revenant au cas $\M$ ponctuée sans plus, \textsc{Quillen} définit, pour $a$, $b$ dans $\M$,
\[
\pi_1(a,b) \;=\; \pi_1(a,b\,;0_a,0_b)
\]
comme bifoncteur sur $(\M\Sigma^{-1})^o\times(\M\Sigma^{-1})$ (avec les notations de 1$^o$)). Si ce que j'avance sans preuve dans 1$^o$) est vrai, il doit s'ensuivre qu'on a un isomorphisme canonique
\[
\pi_1(a,b) \;\simeq\; \pi_1(\Ch^{\Sigma}_{\infty}(a,b)),
\]
où $\Ch^{\Sigma}_{\infty}(a,b)$ est vu comme objet de $\Hot_{\bullet}$, donc son $\pi_1$ est défini. On a en tous cas, sans aucune référence à des cofibrations et des fibrations, un bifoncteur à valeurs dans Groupes, composé dans
\[
(\M\Sigma^{-1})^o\times(\M\Sigma^{-1}) \xrightarrow{\ \ (a,b)\longmapsto \Ch^{\Sigma}_{\infty}(a,b)\ \ } \Hot_{\bullet} \xrightarrow{\ \pi_1\ } (\text{Groupes}).
\]
Si pour $a$ fixé, $b$ variable,
\[
b \;\longmapsto\; \pi_1(a,b)
\]
est représentable, on note l'objet qui le représente par $\Sigma(a)$ (\textit{suspension} -- ne pas confondre ce $\Sigma$ avec $\Sigma \subseteq \mathrm{Fl}(\M)$, qu'il vaut mieux maintenant noter $W$, comme `weak equivalence'). Donc

\msp{page 22}

\[
\mathrm{Hom}(\Sigma a,b) \;\simeq\; \pi_1(a,b),
\]
et $\Sigma a$ dépend fonctoriellement de $a$, sous réserve d'existence. S'il existe toujours, on a
\[
\Sigma :\ \M W^{-1} \;\longrightarrow\; \M W^{-1}.
\]
On définit de même, sous réserve d'existence,
\[
\Omega :\ \M W^{-1} \;\longrightarrow\; \M W^{-1}
\]
par
\[
\mathrm{Hom}(a,\Omega b) \;\simeq\; \pi_1(a,b),
\]
et si tous les deux existent, on a la formule d'adjonction
\[
\mathrm{Hom}(\Sigma a,b) \;\simeq\; \mathrm{Hom}(a,\Omega b) \qquad (\simeq \pi_1\Homext_{\Hot_{\bullet}}(a,b)).
\]

\dgrblock{dgrb438}{Representability and its higher counterpart}
\begin{dcomment}{exposition}{p22:et si tous les deux existent}{Representability and its higher counterpart}
The representability argument is complete once the pointed mapping bifunctor has been constructed. Let $\mathcal C=\mathcal M[W^{-1}]$ and write $F(a,b)=\pi_1\operatorname{Map}^{h}(a,b)$ at zero. If $F(a,-)$ is represented by $\Sigma a$, then a morphism $a\to a'$ induces $F(a',-)\to F(a,-)$; Yoneda gives $\Sigma a\to\Sigma a'$. Uniqueness in Yoneda proves the functor laws. The same argument in the other variable constructs $\Omega$, and the two representing bijections give $\Sigma\dashv\Omega$. The appearance of $0_a,0_b$ in the source's page-21 expression must be read as two copies of $0_{a,b}$, since both endpoints of a path lie in the same mapping space.

Representability of $\pi_1$ alone supplies only a statement on components of loop spaces. To obtain the displayed equivalences of whole mapping spaces one requires homotopy limits and colimits in the enhanced category. In a pointed category with derived mapping spaces and a homotopy pullback defining $\Omega b=0\times_b^h0$, the defining universal property gives
\[
 \operatorname{Map}^{h}(a,\Omega b)
  \simeq *\times_{\operatorname{Map}^{h}(a,b)}^h*
  =\Omega\operatorname{Map}^{h}(a,b).
\]
Dually, the homotopy pushout $\Sigma a=0\amalg_a^h0$ gives the other equivalence. Both respect loop multiplication because concatenation is induced by the same pointed path construction. This is the additional structure behind the source's test~(c), rather than a consequence of an arbitrary bijection on $\pi_0$.

Quillen proves the group-valued bifunctor and its two representabilities for a pointed model category in Chapter I, \S2; the checked retypesetting labels the result Theorem 1.2 \citep{Quillen1967}. The same construction is built into pointed derivators, where the coherent square, rather than a square in the underlying homotopy category, defines suspension and loops \citep[Definition 3.16 and Proposition 3.17]{Groth2013}. The unit being invertible makes $\Sigma$ fully faithful; the counit being invertible makes $\Omega$ fully faithful. If both are invertible, they form an equivalence. Requiring both is the stability condition, not a prerequisite for the existence of the two functors in a pointed model theory.

The following paragraph also uses existence of the appropriate finite products and coproducts. With those in place, a natural multiplication on $\operatorname{Hom}(-,\Omega b)$ is, by Yoneda, a map $\Omega b\times\Omega b\to\Omega b$; the group equations follow by the same argument. Dually $\Sigma a$ is a cogroup. In the enhanced setting, a derived mapping-space functor carries homotopy products in its second variable, and homotopy coproducts in its first, to products of spaces. These are the products and coproducts that justify the source's compatibility with group objects. Existence of ordinary products in the localized category, without a comparison to the enhanced products, would not by itself prove the mapping-space statement. On a pointed derivator $\Omega$ is right adjoint to $\Sigma$ as a morphism of derivators, so it preserves homotopy right Kan extensions, and $\Sigma$ preserves homotopy left Kan extensions; this supplies the left and right exactness anticipated in footnote 8.
\end{dcomment}


% ---- from Der-ChIX-678_p016-018.tex
\printedpage{14}
% head: A. Grothendieck : Dérivateurs, Ch. IX, Édition des Universités de Montpellier II et Paris VII
% foot: Version 12 mars 2002

L'existence de
\[
\begin{aligned}
\Omega :\ & \mathcal{M}W^{-1} \ \longrightarrow\ \mathcal{M}W^{-1}, \quad \text{ou de}\\
\Sigma :\ & \mathcal{M}W^{-1} \ \longrightarrow\ \mathcal{M}W^{-1},
\end{aligned}
\]
voire même des deux à la fois, apparaît comme une hypothèse \cercle{a} extrêmement forte sur $(\mathcal{M}, W)$. Cette condition \textit{plus} celle \cercle{b} que les $\Chu_{\infty}(a,b)$ soient homotopes à des catégories de Thomason-Kan, ou (ce qui revient au même, à peu de choses près?) que $Y \longmapsto \mathrm{Hom}_{\mathcal{M}(Y)W_Y^{-1}}(a_Y, b_Y)$ transforme Hot-équivalences de Cat en bijections (ne serait-ce que pour $Y$ fixé \dots), semble le test crucial pour savoir si `$(\mathcal{M}, \Sigma)$ \texttt{[plutôt} $(\mathcal{M}, W)$\texttt{]} se prête à une théorie homotopique `dans $\mathcal{M}$' '. Je doute que

\msp{page 23}

ces conditions puissent être satisfaites pour $\mathcal{M}$ finie, sauf dans le cas trivial où $\mathcal{M} = \emptyset$, où $\mathcal{M}\Sigma^{-1}$ \texttt{[plutôt} $\mathcal{M}W^{-1}$\texttt{]} est équivalente à la catégorie ponctuelle.

\dgrblock{dgrb439}{A finite test of the proposed axioms}
\begin{dcomment}{development}{p23:ces conditions puissent être satisfaites}{A finite test of the proposed axioms}
Grothendieck suspects that his tests already force an infinite supply of objects and operations. This would be natural for a model of a nontrivial stable homotopy theory, but the tests actually written permit a finite pointed category. The following calculation isolates the missing requirement. We checked the cited localization and model-category sources and the commentary on the earlier chapters; we did not find this particular example discussed as a response to this passage. Its role is diagnostic, without a claim of priority.

\textbf{A finite relative category satisfying tests (a) and (b).}\ours
Let $\mathcal E$ have objects $0,x$, with $0$ a zero object, $\operatorname{End}(x)=\{1_x,z\}$, and $z$ the composite $x\to0\to x$, so $z^2=z\ne1_x$. All other Hom sets are singletons. This is the full subcategory of pointed sets on a singleton and a two-element pointed set. Take $W$ to consist of identities. Then $\mathcal E[W^{-1}]=\mathcal E$, which is neither empty nor equivalent to a point.

We first compute the actual path category. With all backward edges identities, a bilateral eventually constant path is an eventually constant forward chain from $a$ to $b$, with duplicated vertices at the backward steps. Composition of its finite nonidentity part defines a functor
\[
 q:\Ch_\infty^W(a,b)\longrightarrow
       \operatorname{Hom}_{\mathcal E}(a,b)_{\mathrm{disc}}.
\]
Naturality with fixed endpoints makes the composite invariant under morphisms of paths. Choose the section $s$ that inserts a morphism at the forward edge $-1\to0$, with identities elsewhere. Define a third functor $R$ by replacing the part strictly before $0$ by the constant value $a$, retaining the part from $0$ onward, and using the composite $a\to a_0$ at the joining edge. The eventual identity tails make that composite well-defined. There is a natural transformation $R\to\mathrm{id}$: its early components are the composites $a\to a_i$, and its later components are identities. There is also $R\to sq$: its early components are identities and its later components are the composites $a_i\to b$. All the squares commute by composition of forward chains; duplicated backward vertices have identical components. Thus
\[
 \mathrm{id}\longleftarrow R\longrightarrow sq,
 \qquad qs=\mathrm{id}.
\]
This proves a categorical homotopy equivalence to the discrete Hom category. In particular every component is contractible, every $\pi_1$ is zero, and the path categories are homotopy equivalent to discrete categories, which are Thomason--Kan fibrant. This verifies the first formulation of test~(b).

The second formulation is equally direct. Since $W_Y$ contains only identities, no localization changes $\mathcal E(Y)$. A natural transformation between the constant diagrams $a_Y,b_Y$ assigns one arrow $a\to b$ to each connected component of $Y$, because naturality along every arrow forces equality. Hence
\[
 \operatorname{Hom}_{\mathcal E(Y)}(a_Y,b_Y)
 \cong\operatorname{Hom}_{\mathcal E}(a,b)^{\pi_0Y}.
\]
Every Hot-equivalence induces a bijection here. Finally $\Sigma a=0=\Omega a$ represents the identically trivial group-valued bifunctor $\pi_1(a,b)$ in both variables. Test~(a) holds, as do the mapping-space equivalences of test~(c): all their terms are points. The nontrivial endomorphisms of $x$ survive in $\pi_0$, so this is not the terminal exception.

The finite-limit hypothesis introduced on the next page would rule this example out for an elementary reason. If a category has finitely many arrows, say $N$, and all finite powers, then two distinct maps $a\to b$ would give at least $2^n$ maps $a\to b^n$ for every $n$, contradicting $2^n>N$. Thus every Hom set has at most one element. If the category is pointed, the zero maps make every Hom set nonempty, so it is equivalent to the terminal category. This is the finite form of the usual small-completeness argument; see \citet[Exercise 6.3.23(d)]{Leinster2014}. The distinction is therefore exact: the stated tests do not imply the finiteness obstruction, while a finite pointed category with finite products is already trivial before localization.
\end{dcomment}


Un troisième test \cercle{c} qui me paraît important (mais peut-être ça résulte du reste), c'est les formules dans $\Hot_\bullet$
\[
\left\{
\begin{aligned}
\Homext_\bullet(\Sigma a, b) \ &\simeq\ \Homext_\bullet(a, \Omega b)\\
&\simeq\ \Omega\Homext_\bullet(a, b),
\end{aligned}
\right.
\tag*{$(*)$}
\]
où le $\Homext_\bullet$ désigne le Hom externe à valeurs dans $\Hot_\bullet$. Il faut le voir séparément (sous condition d'existence soit de $\Omega$, soit de $\Sigma$) comme deux formules `duales' l'une de l'autre,
\[
\left\{
\begin{aligned}
\Homext_\bullet(a, \Omega b) \ &\simeq\ \Omega\Homext_\bullet(a, b)\\
\Homext_\bullet(\Sigma a, b) \ &\simeq\ \Omega\Homext_\bullet(a, b),
\end{aligned}
\right.
\tag*{$(**)$}
\]
isomorphismes de \textit{groupes} dans $\Hot_\bullet$. (\textbf{NB} $\Omega b$ est un objet groupe de $\mathcal{M}W^{-1}$, car il représente un contrafoncteur à valeurs dans Gr, et de même $\Sigma a$ est un cogroupe de $\mathcal{M}W^{-1}$. Le foncteur $y \longmapsto \Homext(a, y)$ doit commuter aux produits finis, et donc transformer objets groupes en itou. Et de même $x \longmapsto \Homext(x, b)$ doit transformer sommes finies en produits finis, donc cogroupes en groupes (\footnote[8]{Le foncteur $\Omega$ doit provenir d'un endomorphisme du dérivateur $\mathbf{D} = \mathbf{D}_{\mathcal{M},W}$, et celui-ci doit être exact à gauche, et dualement pour $\Sigma$.}).)

Je me suis convaincu que les isomorphismes $(**)$ n'ont de chance d'exister que si on a la condition préalable \cercle{b}.

\newpage

\printedpage{15}
% head: A. Grothendieck : Dérivateurs, Ch. IX, Édition des Universités de Montpellier II et Paris VII
% foot: Version 12 mars 2002

\msp{page 24}

\section*{3\quad Essai de construction du foncteur $\Omega$}

Il n'y a sans doute guère de chance de prouver quelque chose sur le foncteur
\[
\Omega :\ \mathcal{M}W^{-1} \longrightarrow \mathcal{M}W^{-1},
\]
sans tenir en mains une \textit{construction} plus ou moins explicite de celui-ci. Je vais essayer, en m'inspirant de \textsc{Quillen}, de dégager des conditions générales sur $(\mathcal{M}, W)$ (plus générales que celles de \textsc{Quillen}) qui assurent l'existence de $\Omega$, et permettent de le construire. L'idée géométrique remonte à \textsc{Serre-Cartan}. Il s'agit, pour un objet $S$ de $\mathcal{M}$, de construire un objet $\mathrm{Ch}(S)$ \textit{dans $\mathcal{M}$,} jouant le rôle d'un `espace de chemins', et qui doit s'inclure dans un diagramme de factorisation de $\mathrm{diag}_S : S \longrightarrow S \times S$ (on suppose que dans $\mathcal{M}$ les $\varprojlim$ finies, donc les produits finis, existent)
\[
\begin{tikzcd}
S \arrow[d, "i_S"] \arrow[dd, bend right=40, "\mathrm{diag}_S"'] & \\
\mathrm{Ch}(S) \arrow[d, "p_S"] & \\
S \times S & e. \arrow[l]
\end{tikzcd}
\]
La propriété essentielle, c'est que $p$ soit une \textit{`fibration'} (dans un sens convenable, à dégager).

\dgrblock{dgrb440}{Which factorization computes loops}
\begin{dcomment}{exposition}{p24:La propriété essentielle}{Which factorization computes loops}
Grothendieck returns to an actual path object because representability must eventually be proved by a construction. In a pointed model category, choose a fibrant replacement of $S$ and a factorization
\[
 S\xrightarrow{\sim}P S\xrightarrow{(d_0,d_1)}S\times S
\]
of its diagonal, with the first map a weak equivalence and the second a fibration. The fibre over $(0,0)$ computes $\Omega S$. Fibrancy is part of the useful hypothesis: it makes the endpoint maps fibrations to a fibrant object, so their pullbacks compute the intended homotopy fibres. The first map may be chosen to be a trivial cofibration, but its weak-equivalence property is essential. At this point the source is still seeking its axioms, so this is a completion of the proposed construction, not a correction of a theorem claimed to have been proved.

Merely asking for the second map to be a fibration, even adding a cofibration for the first, would be insufficient. In pointed simplicial sets take a nontrivial Kan complex $S=K(A,1)$, let $P S=S\times S$, and let the second map be the identity. The diagonal is a monomorphism, hence a cofibration, and $S$ is fibrant. Nevertheless the indicated fibre is a point, whereas the true loop space has components $A$. What fails is precisely the omitted weak equivalence $S\to P S$.

The representation can be read directly from mapping spaces. In a pointed simplicial model category, let $T$ be cofibrant. Applying $\operatorname{Map}(T,-)$ to the displayed factorization gives a path object for the Kan complex $K=\operatorname{Map}(T,S)$: its endpoint map is a fibration by the simplicial model axiom and its first map is a weak equivalence between fibrant objects. Mapping out of $T$ preserves the pullback that defines the fibre. Therefore
\[
 \operatorname{Map}(T,\Omega S)\simeq\Omega K,
 \qquad
 [T,\Omega S]\cong\pi_1\operatorname{Map}(T,S).
\]
The latter bijection is natural in both variables after replacements. It is independent of the chosen path object because all such objects compute the same homotopy pullback $0\times_S^h0$. For arbitrary model categories the corresponding statement is Quillen's path-object argument, using homotopy classes of left and right homotopies; the checked version gives it as equation (5) in the proof of Theorem 1.2 \citep[Chapter I, \S2]{Quillen1967}.

To replace the mapping space by $\Ch_\infty^W(T,S)$ in the final display of the manuscript requires its comparison with the derived mapping space. A path object in $\mathcal M$ and an interval of zigzags in $\mathrm{Cat}$ play complementary roles; they are not interchangeable definitions. The audit leaves the specific comparison visibly conditional while giving the completed argument in the model setting.
\end{dcomment}


\dgrblock{dgrb441}{Path objects before the full model structure}
\begin{dcomment}{context}{p24:sans tenir en mains une}{Path objects before the full model structure}
Grothendieck is looking for precisely the axioms needed to construct loops. This route had already been isolated by \citet[\S1, pp.~420--421]{Brown1973}, whose categories of fibrant objects retain weak equivalences, fibrations, and path objects without specifying cofibrations. Brown requires pullbacks of fibrations and of trivial fibrations to exist and remain of the same kind; every object is fibrant. His path-object axiom factors the diagonal by a weak equivalence followed by a fibration. These are sufficient structural data for the factorization lemma proved there, which replaces an arbitrary map by a weak equivalence followed by a fibration.

The comparison clarifies the two tentative choices on the next manuscript page. Fibrancy of $S$ supplies part of the working setting, while the essential property of $S\to\mathrm{Ch}(S)$ is its being a weak equivalence. A cofibration condition on that arrow is not needed for Brown's axioms. We use this established precedent to specify a sufficient framework, without suggesting that the manuscript cites Brown or that these axioms alone identify its particular bilateral path categories with derived mapping spaces.
\end{dcomment}


\msp{page 25}

À titre accessoire il peut être utile (peut-être) que $i$ soit une `cofibration' (\footnote[9]{et peut-être que $S \longrightarrow e$ soit également une fibration ?}). On définit $\Omega(S)$ comme la fibre de $\mathrm{Ch}(S)$ sur $S{\times}S$ (par rapport à l'unique morphisme $e = e_{\mathcal{M}} \longrightarrow S\times S$, où $e_{\mathcal{M}}$ est l'objet nul de $\mathcal{M}$).
\[
\begin{tikzcd}[column sep=large]
\mathrm{Ch}(S) \arrow[d] & \Omega S \arrow[l] \arrow[d]\\
S \times S & e. \arrow[l]
\end{tikzcd}
\]

\newpage

\printedpage{16}
% head: A. Grothendieck : Dérivateurs, Ch. IX, Édition des Universités de Montpellier II et Paris VII
% foot: Version 12 mars 2002

C'est un objet de $\mathcal{M}$, qu'on regardera surtout comme un objet de $\mathcal{M}W^{-1}$. Il faut s'assurer qu'on a une bijection fonctorielle en $T$ dans $\mathcal{M}W^{-1}$
\[
\underbrace{\pi_1\Chu^W_\infty(T, S)}_{\overset{\text{déf}}{=}\,\pi_1(T,S)} \ \simeq\ \underbrace{\mathrm{Hom}_{\mathcal{M}W^{-1}}(T, \Omega S)}_{\overset{\text{déf}}{=}\,\pi_0\Chu^W_\infty(T,\Omega S)}.
\]
Pour cela, il s'impose de comparer
\[
\Chu^W_\infty(T, \Omega S) \qquad\quad \text{et} \qquad\quad \underbrace{\Omega_{\Hot_\bullet}(\Chu^W_\infty(T, S))}_{\overset{\substack{\text{`par définition'}\\ \text{pour $\Omega$ de $\mathrm{Hot}_\bullet$}}}{=}\ \Chu_\infty(\mathrm{Ch}^W_\infty(T,S);0,0)} \qquad ,
\]
puisqu'il s'agit de construire une bijection entre leurs $\pi_0$.

\dgrblock{dgrb442}{The loop object of a complex, with its contraction}
\begin{dcomment}{examples}{p25:puisqu'il s'agit de construire une bijection}{The loop object of a complex, with its contraction}
For complexes the proposed construction can be written down without any choices, and its fibre displays the sign in the loop shift. Use cochain complexes of $k$-modules with the model structure whose weak equivalences are quasi-isomorphisms and whose fibrations are degreewise surjections. Every complex is fibrant. For a complex $B$, define
\[
 (P B)^n=B^n\oplus B^n\oplus B^{n-1},\qquad
 d(x,y,h)=(dx,dy,y-x-dh).
\]
The third component of $d^2$ is $dy-dx-d(y-x-dh)=0$, so this is a complex. The maps
\[
 i(b)=(b,b,0),\qquad p(x,y,h)=(x,y),\qquad r(x,y,h)=x
\]
are chain maps, $pi$ is the diagonal, and $p$ is a degreewise surjection. To verify the required weak equivalence, define the degree $-1$ homotopy $H(x,y,h)=(0,h,0)$. Direct calculation gives
\[
 (dH+Hd)(x,y,h)=(0,y-x,h)
               =(\mathrm{id}-ir)(x,y,h),\qquad ri=\mathrm{id}.
\]
Thus $i$ is a chain homotopy equivalence. Its fibre over zero under $p$ is the complex with terms $B^{n-1}$ and differential $-d$, namely $B[-1]$. We have obtained $\Omega B=B[-1]$ and, dually, $\Sigma B=B[1]$.

For a cofibrant replacement $Q T\to T$, the corresponding Hom complex satisfies
\[
 H^0\operatorname{Hom}^\bullet(QT,B[-1])
   =H^{-1}\operatorname{Hom}^\bullet(QT,B)
   =\pi_1\operatorname{Map}^{h}(T,B),
\]
which verifies the source's required bijection in this example. More than the bijection is visible: applying the good nonpositive truncation and Dold--Kan to these Hom complexes identifies the whole mapping space into $B[-1]$ with the loop of the mapping space into $B$. All maps used above commute with chain maps $B\to B'$, so the construction is functorial before localization. This is exactly the kind of explicit factorization that the closing paragraph asks us to find.
\end{dcomment}
\dgrchapter{X}{Chapter X}{\texorpdfstring{Comparaison de Cat avec $\boldsymbol{\Delta}^{\wedge}$}{Comparaison de Cat avec Δ∧}}
\dgrpart{0}{VIII-X}{Manuscript pages 1--5}{VIII-X\enspace\textperiodcentered\enspace Chapter X\enspace\textperiodcentered\enspace manuscript pages 1--5}{Chapter X \textperiodcentered\ \textup{VIII-X}}
% ---- from Der-ChX-789_p001-005.tex
% ===================== page 1 (p001) =====================
\printedpage{}
% no page number printed (title page)
% head: (none)

\begin{center}
{\LARGE \textsf{LES DÉRIVATEURS}}

\bigskip

{\large \textbf{ALEXANDRE GROTHENDIECK}}

\vspace{3em}

{\large \textbf{Chapitre X}}

\medskip

{\large \textbf{Comparaison de} Cat \textbf{avec} $\bDelta^{\wedge}$}
\end{center}

\clearpage

% ===================== page 2 (p002) =====================
\printedpage{}
% no page number printed (editorial note)
% head: (none)

Ce texte a été déchiffré et transcrit en \LaTeXe{} par M. Künzer. Il a été édité par M. Künzer, J. Malgoire, et G. Maltsiniotis. La transcription est aussi fidèle que possible au manuscrit. Pour les quelques corrections évidentes, ou rares commentaires des éditeurs, ainsi que pour la numérotation originale des pages du manuscrit, les caractères de machine à écrire \texttt{[typewriter]} entre crochets sont utilisés. Un point d’interrogation entre crochets signifie que l’on n’est pas sûr du mot qui précède.

Cette édition est provisoire. Les remarques, commentaires et corrections sont bienvenus. Envoyer un message à :

\textbf{maltsin@math.jussieu.fr}

\hfill G. Maltsiniotis\hspace*{3em}

\clearpage

% ===================== page 3 (p003), printed 1 =====================
\printedpage{1}
% head: A. Grothendieck : Dérivateurs, Ch. X, Édition des Universités de Montpellier II et Paris VII
% foot: Version 9 avril 2002

\msp{page 1}

\section*{1\quad Rappel des notions de D-asphéricité, et lien avec le ‘localisateur fondamental $W_{\mathbf{D}}$’}

Pour un dérivateur $\D$, il s’agit des notions suivantes (cf. I) :

\dgrblock{dgrb443}{The homotopy shapes seen by a coefficient theory}
\begin{dcomment}{exposition}{p1:il s’agit des notions suivantes}{The homotopy shapes seen by a coefficient theory}
Grothendieck now returns to the part of a derivator that can be expressed without choosing a particular category of coefficients. The relevant equivalences concern \emph{constant} coefficients. Write $p_A:A\to e$ and $K_A=p_A^*K$. Then $u:A\to B$ belongs to $W_{\D}$ precisely when, for every $K,L\in\D(e)$, restriction gives a bijection
\[
 \Hom_{\D(B)}(K_B,L_B)\ \longrightarrow\
 \Hom_{\D(A)}(K_A,L_A).
\]
This recalls the distinction explained in the commentary on Chapter~I, manuscript p.~61 (``Equivalence measured by constant coefficients''). It is not the assertion that $u^*:\D(B)\to\D(A)$ is an equivalence on arbitrary coherent diagrams. The distinction is what makes $W_{\D}$ a theory of homotopy \emph{shapes}, while $\D$ still records the coefficients carried by those shapes.

The two comma categories retain different orientations. An object of $A/b$ is an arrow $u(a)\to b$; an object of $b\backslash A$ is an arrow $b\to u(a)$. Accordingly, the fibre of $A/a\to B/u(a)$ over $\beta:b\to u(a)$ consists of factorizations of $\beta$ by arrows $a'\to a$ lying over it. Its asphericity is the lifting criterion for smoothness, with the opposite criterion for properness. These are the categorical forms of the two base-change conditions, not assertions that the strict fibres $A_b$ alone control base change. The later formulation is \citet[Theorem~1.2 and its dual]{Maltsiniotis2005Asphericite}.
\end{dcomment}


\begin{itemize}
\item[(i)] \textit{Objet} $\D$-\textit{asphérique} de Cat (i.e. $Z \longrightarrow e$ est une $\D$-équivalence). NB Je préfère dire ‘$\D$-asphérique’ au lieu de ‘$\D$-contractile’, qui prête parfois à confusion.

\item[\textcircled{ii}] $\D$-équivalence pour une flèche $f : X \longrightarrow Y$ dans Cat.

\item[(iii)] \textit{Foncteur} $f : X \longrightarrow Y$ $\D$-\textit{asphérique} (i.e. les $X/y$ sont $\D$-asphériques) et $\D$-\textit{coasphérique} (les $y\backslash X$ sont $\D$-asphériques). C’est ce que j’avais appelé ‘$\D$-équivalence stricte’ (ou ‘costricte’), ou foncteur homologiquement $\D$-cofinal (ou cohomologiquement $\D$-cofinal), mais la terminologie ‘asphérique’ me semble la plus simple et la meilleure. (Pêchée dans PS, Pursuing Stacks.)

\item[(iv)] \textit{Foncteur} $f : X \longrightarrow Y$ $\D$-\textit{lisse,} $\D$-\textit{propre} (i.e. tel que les fibres des foncteurs induits $X/x \longrightarrow Y/f(x)$, respectivement $x\backslash X \longrightarrow f(x)\backslash Y$, soient $\D$-asphériques).

\item[(v)] $\D$-\textit{fibration} $f : X \longrightarrow Y$ (notion autoduale), i.e. telle que pour $Y'' \xrightarrow{u} Y'$ au dessus de $Y$, si $Y''$, $Y'$ \texttt{[sont]} Hot-asphériques, alors $u_X : X'' \longrightarrow X'$ dans
\[
\begin{tikzcd}
X \arrow[d] & X' \arrow[l] \arrow[d] & X'' \arrow[l] \arrow[d] \\
Y & Y' \arrow[l] & Y'' \arrow[l]
\end{tikzcd}
\]
est une $\D$-équivalence (\footnote{On dira aussi que $f$ est $\D$-\textit{fibrant.} (\textbf{NB} Il n’y a pas de notion ‘$\D$-cofibrant’ à ma connaissance.)}). La $\D$-\textit{fibration} est dite \textit{triviale} (ou \textit{asphérique}) si et seulement si ses fibres sont asphériques, ou encore si et seulement si $f$ est asphérique, ou si et seulement si $f$ est coaspérique. Cela signifie aussi que $f$ est une $\D$-équivalence \textit{universelle,} % [sic] coaspérique -> coasphérique, letter h missing
\msp{page 2}
i.e. que pour tout $Y' \longrightarrow Y$, $X' = X \times_Y Y' \longrightarrow Y'$ est une $\D$-équivalence.

\item[(vi)] On a aussi des variantes affaiblies de la notion de $\D$-fibrations, savoir les ‘$\D$-\textit{fibrations à gauche’,} ou les $f$ ‘\textit{cohomologiquement} $\D$-\textit{fibrants’,} et de même les ‘$\D$-\textit{fibrations à droite’,} i.e. les $f$ ‘\textit{homologiquement} $\D$-\textit{fibrants’} (\footnote{Nouvelle terminologie : fibrations cohomologiques (resp. homologiques) \textit{faibles,} cf. I, pages 111-126 pour une révision corrigée des diverses notions de ‘$\D$-fibrations’.}). La notion ‘gauche’
\printedpage{2}
% head: A. Grothendieck : Dérivateurs, Ch. X, Édition des Universités de Montpellier II et Paris VII
% foot: Version 9 avril 2002
ou cohomologique pour $f : X \longrightarrow Y$ signifie que les foncteurs $X_y \hookrightarrow X/y$ ($y \in \Ob Y$) et $X/y \longrightarrow X/y'$ (pour $(u : y \longrightarrow y') \in \Fl Y$) sont des $\D$-équivalences. Et dualement avec les $X_y \hookrightarrow y\backslash X$, $y'\backslash X \longrightarrow y\backslash X$ pour la notion ‘droite’ ou ‘homologique’. Les $\D$-fibrations ne sont autres que les $\D$-fibrations gauches (cohomologiques) universelles ou les $\D$-fibrations droites (homologiques) universelles.
\end{itemize}

\dgrblock{dgrb444}{Why triviality is a universal condition}
\begin{dcomment}{exposition}{p2:fibrations gauches (cohomologiques) universelles}{Why triviality is a universal condition}
The point of the fibration condition is to make the fibre calculations survive a change of base. The qualifier ``universal'' in the next sentence therefore carries real mathematical content; the earlier distinctions are reviewed in the commentary on Chapter~I, manuscript p.~104 (``The universal criterion, with the quantifiers restored'').

Here is a useful direct verification of the asserted triviality criterion. Suppose $f:X\to Y$ is a $\D$-fibration and every $X_y$ is $W_{\D}$-aspherical. After any change of base $B\to Y$, the same hypotheses hold. For $b\in B$, apply the definition to $e\to B/b$, selecting its final object. Both bases are ordinarily contractible, so
\[
 (X\times_Y B)_b\longrightarrow (X\times_Y B)/b
\]
lies in $W_{\D}$. Hence every category on the right is $W_{\D}$-aspherical. The asphericity theorem gives $X\times_Y B\to B\in W_{\D}$. Conversely, a universal $W_{\D}$-equivalence has aspherical fibres by base change to points. It is a $\D$-fibration: in its defining square the two projections to $Y''$ and $Y'$ lie in $W_{\D}$, and $Y''\to Y'$ does also because both bases are ordinarily contractible; two-out-of-three gives the required arrow between pullbacks. Applying the same comparison with $Y/y$ and $y\backslash Y$ proves the two equivalent comma-category criteria. The ordinary contractibility of the test bases, specified in the manuscript, has not been replaced by $W_{\D}$-asphericity.
\end{dcomment}


Finalement, toute cette kyrielle de notions se déduit de la seule notion de $\D$-équivalence (ii). Si je désigne par
\[
W_{\D} \subseteq \Fl(\Cat)
\]
l’ensemble de ces $\D$-équivalences, on trouve ce que dans Pursuing Stacks j’ai appelé un ‘localisateur’ (localiseur) dans Cat (\footnote{Mais attention, je suppose ici que le domaine Diag de $\D$ est Cat tout entier!}). Dans Pursuing Stacks c’était une telle partie
\[
W \subseteq \Fl(\Cat),
\]
et non la donnée d’un dérivateur, qui
\msp{page 3}
était pris comme la donnée de départ fondamental pour faire de l’algèbre homotopique en tous genres. J’avais d’ailleurs développé une kyrielle de notions du même genre (i) à (vi), en partant de cette donnée $W$. Il se pourrait que les énoncés pour les notions que j’ai développé dans I, en termes d’un dérivateur, soient valables plus généralement en termes d’un localisateur $W$, dans la mesure du moins où ils gardent un sens. (Donc pour l’instant je mets à part les énoncés proprement cohomologiques ou homologiques, qui font directement appel à $\D$.) Encore qu’il se pourrait que moyennant des conditions à expliciter sur $W$, qui impliquent que $(\Cat, W)$ définissent bel et bien un dérivateur $\D_{\Cat,W} = \D$, on arrive encore à interpréter $W$ comme un $W_{\D}$, de sorte que les énoncés cohomologiques et homologiques gardent un sens.

\dgrblock{dgrb445}{Realizing a localizer by a derivator}
\begin{dcomment}{later}{p3:les énoncés cohomologiques et homologiques gardent un sens.}{Realizing a localizer by a derivator}
Grothendieck asks for a coefficient theory realizing the proposed weak equivalences. A precise answer is available under an accessibility hypothesis. Fix a universe for small diagram shapes and a larger universe for coefficient categories, and call a fundamental localizer accessible when it is generated by a small set of functors. \citet[Theorem~5.2.15]{Cisinski2006} constructs a cofibrantly generated, left proper model structure on $\Cat$ with weak equivalences $W$. The model-category construction of derivators then gives
\[
 \D_W(A)=\Cat^{A^o}\bigl[(W^{A^o})^{-1}\bigr]
\]
with both homotopy Kan extensions \citep[Theorem~6.11 and its dual]{Cisinski2003}. The comparison with simplicial presheaves and the recovery of $W$ from constant-coefficient tests are supplied by \citet[Lemma~4.2.13, Theorem~4.2.15, and Corollary~4.2.18]{Cisinski2006}. Thus $W_{\D_W}=W$. Accessibility is the hypothesis of this realization statement; it is not implicit in the definition of an arbitrary derivator.

Realization does not determine a pre-existing coefficient theory. For a complete and cocomplete category $\mathcal C$, the represented derivator satisfies
\[
 \Hom_{\mathcal C^{A^o}}(K_A,L_A)
   =\operatorname{Map}\bigl(\pi_0 A,\Hom_{\mathcal C}(K,L)\bigr).
\]
Indeed, naturality says exactly that the component of a natural transformation is constant along every zigzag in $A$. For $\mathcal C=\mathrm{Set}$ or $\mathrm{Ab}$, a Hom set with at least two elements detects whether $\pi_0(u)$ is bijective: a missing component prevents injectivity of precomposition, and two identified components prevent surjectivity. Both derivators consequently give the same localizer $W_0$ of component bijections. Their categories of coefficients are not equivalent, since $\mathrm{Ab}$ has a zero object and $\mathrm{Set}$ does not. The equality $W_{\D_W}=W$ answers the question on these pages without identifying the original $\D$ with $\D_W$.
\end{dcomment}


J’ai envie d’abord de regarder dans Pursuing Stacks les axiomes que j’ai dégagé pour les localisateurs, et vérifier s’ils sont tous vérifiés pour
\msp{page 4}
$W_{\D}$. Ces conditions L 1 à L 5 sont énoncées dans Pursuing Stacks, §64 (page 164).

\textbf{L 1}\enspace Saturation de $W_D$. O.K.

Cela résulte formellement du fait qu’on a un foncteur
\[
\Cat \xrightarrow{\varphi} C \qquad (= \Cat^o)
\]
tel que pour $u \in \Fl(\Cat)$, $u \in W_{\D} \iff \varphi(u)$ isomorphisme. (Prendre pour $\varphi(X)$ la sous-catégorie de $\D(X)$ formée des $\D$-coefficients \textit{constants} sur $X$).

\dgrblock{dgrb446}{Making the saturation argument literal}
\begin{dcomment}{exposition}{p4:isomorphisme. (Prendre pour}{Making the saturation argument literal}
Grothendieck's argument is sound once the category of constants is represented with its coefficient labels retained. Define $C_{\D}(A)$ to have objects $K\in\D(e)$ and morphisms
\[
 C_{\D}(A)(K,L)=\Hom_{\D(A)}(K_A,L_A).
\]
Restriction is the identity on these object labels. Thus $u\in W_{\D}$ if and only if $C_{\D}(u)$ is an \emph{isomorphism} of categories, whereas the essentially full subcategory of constant objects gives the equivalent formulation using equivalences of categories. This avoids choosing representatives for coincident or isomorphic constant diagrams.

The construction proves strong saturation. If the image of $u$ becomes invertible in $\Cat[W_{\D}^{-1}]$, the factorization of the contravariant functor $C_{\D}$ through that localization makes $C_{\D}(u)$ invertible; hence $u\in W_{\D}$. In \emph{Pursuing Stacks}, \S64, ``saturation'' is the weaker list of identities, two-out-of-three, and the two-sided-composite condition. The present argument proves more than that list requires \citep[\S64, manuscript pp.~163--165]{Grothendieck1983}.
\end{dcomment}


\textbf{L 2}\enspace “$\bDelta^1$ est $W$-asphérique sur $e$”, i.e. pour tout $X$ dans Cat, $X \times \bDelta^1 \longrightarrow X$ est dans $W$.

Cela signifie aussi que $\bDelta^1 \longrightarrow e$ est une $\D$-fibration triviale, ce qui est bien le cas.

\clearpage

% ===================== page 5 (p005), printed 3 =====================
\printedpage{3}
% head: A. Grothendieck : Dérivateurs, Ch. X, Édition des Universités de Montpellier II et Paris VII
% foot: Version 9 avril 2002

\textbf{L 3}\enspace Si $X \in \Ob\Cat$ a un objet final, alors $X \longrightarrow e$ est dans $W$.

En fait, nous exigeons pour les dérivateurs que ceci soit vrai en supposant seulement $X$ asphérique (au sens absolu), et on en conclut même que $X \longrightarrow e$ est $W$-asphérique, i.e. $X \times Y \xrightarrow{\mathrm{pr}_2} Y$ est dans $W$ pour tout $Y$ dans Cat (\footnote{En fait L 2 $\Longrightarrow$ L 3.}).

\textbf{L 4}\enspace Si $f : X \longrightarrow Y$ est tel que $X/y \longrightarrow Y/y$ sont dans $W$, alors $f$ est dans $W$.

C’est OK.

\msp{page 5}
\textbf{L 5}\enspace Si $X$, $Y$ sont Cat-fibrées sur $S$, et $X \longrightarrow Y$ est un $S$-morphisme \textit{cartésien,} induisant sur les fibres des $X_s \longrightarrow Y_s$ qui sont dans $W$, alors $f$ est dans $W$.

Pour $W_{\D}$, c’est la proposition 6, page 71 de I.

Ces axiomes L 1 à L 5 (équivalent à L 1, L 3$'$ ($\bDelta^1 \longrightarrow e$ est dans $W$) et L 5) sont dans la nature d’axiomes de stabilité, assurant que $W$ est assez grand. Il y a une autre condition qui limite $W$ à ne pas être trop grand :

\dgrblock{dgrb447}{Completing the reduction to three axioms}
\begin{dcomment}{exposition}{p5:assurant que $W$ est assez grand.}{Completing the reduction to three axioms}
The shortened axiom list expresses a useful economy: integration of families, together with one interval, already forces the basic homotopy calculus. The reduction is sketched in \emph{Pursuing Stacks}, \S64. We spell it out, using the twisted-arrow construction of Chapter~I rather than its mapping-cone presentation; compare the commentary on Chapter~I, manuscript p.~73 (``The twisted-arrow category and its two projections'').

Assume L~1, L~3$'$, and L~5. Applying L~5 to the constant fibrations $A\times[1]\to A$ gives L~2. Their two sections belong to $W$ by two-out-of-three. A natural transformation $F\Rightarrow G$, viewed as a functor $A\times[1]\to B$, therefore shows that $F\in W$ if and only if $G\in W$. A category with a final or initial object retracts to that object, with its other composite naturally homotopic to the identity. Saturation gives L~3 and its opposite.

To pass between the two comma orientations, let $T(A)$ have objects arrows $a\to b$ and morphisms from $h:a\to b$ to $h':a'\to b'$ pairs
\[
 r:a\to a',\qquad s:b'\to b,\qquad s h'r=h.
\]
Both projections $T(A)\to A$ and $T(A)\to A^o$ are Grothendieck fibrations. Their fibres are $(a\backslash A)^o$ and $A/b$, respectively, and have final objects. L~5 consequently puts both projections in $W$. The construction is natural in $A$, so its two commutative squares and two-out-of-three give
\[
 u\in W\quad\Longleftrightarrow\quad u^o\in W.
\]
This also establishes the final assertion L~d).

Next let $u:A\to B$ have $W$-aspherical categories $b\backslash A$. The Grothendieck construction
\[
 E=\int_{b\in B}(b\backslash A)\longrightarrow B
\]
is a fibration with these fibres, hence lies in $W$ by L~5. Its objects are triples $(b,a,\alpha:b\to u(a))$. The forgetful functor $r:E\to A$ has the section $i(a)=(u(a),a,1)$, satisfies $ri=1$, and has a natural transformation $1_E\Rightarrow ir$ with component $(\alpha,1_a)$. Thus $r$ and $i$ belong to $W$, and $u$ does also. Opposite invariance proves the analogous assertion for $A/b$. Under the hypothesis of L~4, the final object of $B/b$ makes every $A/b$ aspherical by two-out-of-three, so L~4 follows. Conversely, L~1--L~5 include L~1 and L~5 and imply L~3$'$ by L~3. This proves the claimed equivalence of axiom systems. For the modern relative form, suppose $A\to B$ lies over $C$ and every $A/c\to B/c$ is in $W$. Opposite invariance gives objectwise equivalences $(A/c)^o\to(B/c)^o$. These diagrams define fibrations over $C^o$, so L~5 applies to their Grothendieck constructions; their forgetful functors are deformation retractions to $A^o$ and $B^o$. Thus $u^o$, and hence $u$, belongs to $W$. This yields the relative Theorem~A axiom of \citet[Definition~1.1.3]{Cisinski2004Minimal}. No model structure is used.
\end{dcomment}


\textbf{L a)}\enspace Si $f \in W$, alors $\pi_0(f)$ est bijectif.

Ceci n’est pas satisfait pour tous les $W_{\D}$ – p. ex. si $\D$ est le dérivateur \textit{nul,} alors $W_{\D} = \Fl\Cat$! Mais pour les dérivateurs utiles en pratique, cela semble toujours satisfait.

\dgrblock{dgrb448}{A nonzero coefficient theory that forgets components}
\begin{dcomment}{examples}{p5:Mais pour les dérivateurs utiles en pratique}{A nonzero coefficient theory that forgets components}
The additional condition L~a) restricts the amount of homotopy theory the coefficients may forget. Even a nonzero derivator need not satisfy it. Take the complete lattice $\mathbf 2=\{0<1\}$ and its represented derivator
\[
 \D_{\mathbf 2}(A)=\operatorname{Fun}(A^o,\mathbf 2).
\]
For nonempty $A$, the Hom sets between the two constant diagrams are the Hom sets of $\mathbf 2$: all are singletons except $\Hom(1_A,0_A)=\varnothing$. For $A=\varnothing$ there is one diagram and one morphism. Therefore
\[
 u:A\to B\in W_{\D_{\mathbf 2}}
 \quad\Longleftrightarrow\quad
 (A=\varnothing\ \Longleftrightarrow\ B=\varnothing).
\]
In particular the fold $e\amalg e\to e$ belongs to this localizer, although its map on components is not injective. This is the usual coarse localizer $W_{-1}$ \citep[Examples~3.3.5--3.3.6]{Cisinski2006}. The terminal coefficient category gives the still larger localizer of all functors, exactly as Grothendieck observes. By contrast, coefficients with a Hom set of at least two elements detect components in the represented examples calculated above. His remark about useful examples is a comment on the intended coefficients, not an additional derivator axiom.
\end{dcomment}


\textbf{L b), c)}\enspace Le foncteur
\[
\Cat \longrightarrow \Cat\, W^{-1}
\]
commute aux sommes finies et aux produits finis (\footnote{Ça marche si $(\Cat, W)$ provient d’un triple de Quillen (avec des Q-fibrations, Q-cofibrations, Q-équivalences faibles).}).

À examiner ce que ça signifie en termes concrets sur $W$ – je n’ai pas vérifié que c’est O.K. pour les $W_{\D}$.

\dgrblock{dgrb449}{Finite products and sums already descend}
\begin{dcomment}{status}{p5:À examiner ce que ça signifie en termes concrets}{Finite products and sums already descend}
The questions L~b) and L~c) have a formal answer under the axioms already retained. If $u:A\to B$ is in $W$, the map of constant fibrations over any $C$ gives
\[
 u\times1_C:A\times C\longrightarrow B\times C\quad\in W
\]
by L~5. Applying L~5 over a discrete indexing category likewise shows that coproducts of arrows in $W$ belong to $W$. These closure properties are also \citet[Propositions~1.1.5 and 1.1.7]{Cisinski2004Minimal}.

Closure alone should be connected to the universal properties of the claimed products and sums. Put $H=\Cat[W^{-1}]$. The localization of $\Cat\times\Cat$ at the componentwise weak equivalences is $H\times H$: a pair of zigzags can be interleaved by holding the other coordinate fixed, and the commutation squares between the two coordinates remain valid after inversion. The diagonal and the product functor are homotopical, so their adjunction, including its unit and counit, descends to $H$. Thus the image of $A\times B$ represents the product of the images of $A$ and $B$. Descending the adjunction between coproduct and diagonal proves the corresponding coproduct assertion. The nullary adjunctions give the initial and terminal objects as well. Consequently the canonical localization preserves finite sums and finite products. A Quillen structure is unnecessary here; the explicit closure argument supplies the homotopy invariance needed for the \emph{ordinary} products and sums appearing in the question.
\end{dcomment}


\textbf{L d)}\enspace Si $f : X \longrightarrow Y$ dans $W$, alors $f^o : X^o \longrightarrow Y^o$ aussi.

C’est le corollaire 2, page 75 de I.
\begin{otherlanguage}{english}
\dgroverview{VIII-X}{Chapter X \textperiodcentered\ \textup{VIII-X} \textperiodcentered\ Overview}
\selectlanguage{english}
\noindent\textbf{Integration, spaces of morphisms, and the choice of weak equivalences.}

These three chapters investigate how much of homotopy theory can be reconstructed from operations on ordinary categories, provided that one keeps the right forms of coherence. Chapter VIII isolates integration for relative $1$-types. Chapter IX returns to the path categories introduced earlier, first to extract an abstract theory of composition and then to compare it with Quillen's constructions. The opening of Chapter X changes the point of departure: instead of beginning with a derivator and its coefficients, Grothendieck asks whether a class of weak equivalences of categories can itself generate the theory. The movement is from a concrete adjunction, through the structure of its spaces of morphisms, to the axioms that make such constructions possible.

These pages comprise all of Chapters VIII and IX and the opening of Chapter X, twenty-seven pages of the edition. Manuscript references must always include the chapter: they run through VIII.1--9, IX.1--25, and X.1--5. Our audit follows those separate sequences. Original statements, including the edition's retained slips and editorial interventions, remain distinguishable from the commentary. The two doubtful readings are settled against the edition; neither licenses a mathematical change to Grothendieck's text.

\medskip\noindent\textbf{Relative groupoids and a genuine biadjunction.}

The decisive construction of Chapter VIII is simple enough to write down on an object of the base. For $f:S\to T$ and a category $G$ fibred in groupoids over $S$, set
\[
 (f_!G)_t=\Pi(t\backslash_TG).
\]
The comma category records a point of $G$ together with an arrow from $t$ to its image in $T$; groupoid completion turns this record into a homotopy $1$-type. Its use is not limited to an already fibred category. For any $X\to T$, the same construction gives a relative groupoid $R_TX$, and extension along $X\to R_TX$ is universal for functors into categories fibred in groupoids. Cartesian lifts construct the extension, while their uniqueness supplies the coherence.

This universal property completes the adjunction verification left unfinished in the manuscript. It proves an equivalence of Hom-groupoids, not only a bijection on isomorphism classes. Both triangle comparisons arise from the same extension-and-restriction construction. Transitivity of integration is then the uniqueness of the left biadjoint to composite restriction. The point of the $2$-categorical formulation is thus practical: the automorphisms that would have been discarded by taking components are precisely what make the comparison usable in a further composition. The base-change theorem has the corresponding later expression in the theory of asphericity structures and smooth functors \citep[\S5]{Maltsiniotis2005Asphericite}.

The final descent statements require a more careful separation. The proper or smooth hypotheses of Theorem 3 support its assertions about faithfulness and about the essential image described by coefficients that are \emph{constant on each strict fibre}. They do not justify replacing this condition by the assertion that every vertical transport is an equivalence. A locally constant groupoid family can carry degree-two coherence even when its base is simply connected. We give a finite example over the fourteen-element face poset of the boundary of a tetrahedron. Every fibre is $B(\mathbf Z/2)$ and every transport is the identity, but a nonzero normalized $2$-cocycle prevents the total category from admitting a section. Consequently it cannot be equivalent to a constant family. This disproves part 3 and the subsequent corollary as printed, while leaving the preceding fibrewise-constant criterion intact.

The repair is expressed by higher monodromy: equivalence classes of locally constant groupoid families with a fixed fibre are controlled by maps to its automorphism $2$-type. A $1$-connected base removes ordinary monodromy but can retain the degree-two class; $2$-connectivity removes both \citep[Corollary~2.2.4 and Theorem~2.2.5, arXiv version]{PoleselloWaschkies2005}. A separate, shorter counterexample shows that section groupoids alone do not test all the proposed faithfulness conditions. The correct section tests must also be made after passing to the total category of the first coefficient. Both corrections concern the quantifiers and the coherence in descent, rather than the integration construction itself.

The general faithfulness question also admits an exact answer without smoothness or properness. For each arrow of $T$, form its category of factorizations through objects in the image of $f$. Restriction is faithful on transformations, fully faithful on Hom-groupoids, or an equivalence on Hom-groupoids precisely when these factorization categories are respectively nonempty, connected, or $1$-connected. We prove this three-level criterion by the counit on representables and coherent density, alongside the known untruncated homotopical-epimorphism criterion \citep[Theorem~6.10, arXiv version]{GrothStovicek2016Trees}.

\medskip\noindent\textbf{What composition gives, and what it presupposes.}

Chapter IX begins with a strong idea: composition of paths should explain endpoint functoriality. In a monoidal category $\mathcal H$, Hom objects $C(a,b)$ with associative composition and coherent units determine an ordinary category with morphisms $\operatorname{Hom}_{\mathcal H}(I,C(a,b))$. Inserting these elements into composition produces the bifunctor $C_0^{\mathrm{op}}\times C_0\to\mathcal H$. The proof is formal once the hypotheses are in place. Neither symmetry nor an internal Hom is needed. This is the familiar passage from an enriched category to its underlying ordinary category \citep[\S\S1.1--1.3]{Kelly1982}.

There are two limits to what this argument proves in the passage under review. First, the endpoint maps give isomorphisms only when those endpoint maps are invertible. The early corollary omits that hypothesis, and a discrete two-object example disproves the wording; the manuscript itself restores it on page 8. Second, the abstract argument assumes associative composition on the particular objects $\Ch_\infty(a,b)$. It does not construct that composition. Grothendieck sketches the finite-stage argument for the ind-object $\mathrm{CH}_\infty$ and explicitly leaves the corresponding realized-space verification unfinished.

The possible shortcut through finite tests does not solve this difficulty. The earlier definition of $\Hot_0$ uses homotopy types of finite categories, not just finite CW complexes. The resulting restricted Yoneda functor is not conservative, by the general theorem on small families of unpointed tests, and it is not faithful, by nonconcreteness of the homotopy category \citep[Theorem~2.1]{ArlinChristensen2023} and \citet[pp.~1, 5 of the reprint]{Freyd1970}. Thus equality after passing to those presheaves cannot replace a compatible homotopy on the original mapping spaces. The commentary also distinguishes an ind-object, an ordinary colimit in $\Hot$, and a homotopy colimit. An explicit increasing sequence of finite discrete sets has an ordinary colimit without being essentially constant as an ind-object. A finite-category test such as $BC_2$, on the other hand, is not compact with respect to all filtered homotopy colimits. These examples delimit the arguments that the comparison maps actually require.

Modern localization supplies the intended kind of mapping-space theory. The hammock localization has coherent mapping spaces and a component category equal to the ordinary localization \citep[\S2]{DwyerKan1980Calculating}. This identifies the modern object to which Grothendieck's construction should be compared; it does not by itself prove a comparison for his exact bilateral stabilization and transformation conventions. We keep that model-specific verification explicit. It is not described as a new open problem in the general theory of localization.

\medskip\noindent\textbf{Pointed and linear refinements.}

The later part of Chapter IX asks that the zero map become a specified basepoint in every mapping space. A basepoint in the unpointed homotopy category is insufficient: based homotopies remember data lost under free homotopy. The commentary constructs the pointed morphism associated to a map and a chosen path between basepoints, including its composition law. It also explains the significance of the manuscript's two corrective footnotes. The Cartesian unit in pointed spaces has only one map to every object; the unit that recovers all components is $S^0$ for the smash product. Composition must therefore annihilate either zero input. Forgetting basepoints carries this structure to unpointed composition through a lax monoidal comparison.

Some questions at the level of ordinary localization have particularly direct answers. Initial and terminal objects survive localization without extra hypotheses, by naturality of their unique arrows along both ordinary and inverted generators. Suitable closure under finite products or sums lets the corresponding adjunctions descend. For an additive category, stability of the denominator class under finite direct sums already yields an additive ordinary localization \citep[Lemma~2.2.1]{Krause2021}. The stronger request for complexes of morphisms requires an enhanced setting. We treat complexes of modules explicitly, with connective truncations of derived Hom, derived tensor composition, and the Dold--Kan comparison. This makes the grading visible: positive homotopy groups of the mapping space are negative cohomology groups of the Hom complex.

The suspension and loop objects are then representability statements. Representing the fundamental-group bifunctor yields an adjunction; the equivalences of whole mapping spaces arise from homotopy pushouts and pullbacks. An explicit path complex verifies the loop construction with its differential, its contraction, and its fibre. We also develop a finite counterexample to the proposed force of tests (a) and (b): the full category of pointed sets on one and two elements, with only identities inverted, satisfies both tests and even the displayed mapping-space identities. Its loop and suspension functors are zero, while its ordinary category is nontrivial. The missing strength is therefore not finiteness of the category alone. This calculation is marked as the commentators' development, without a priority claim.

\medskip\noindent\textbf{Localizers and the information retained by coefficients.}

The opening of Chapter X extracts a class $W_{\D}$ by testing maps of indexing categories against constant coefficients. It must be distinguished from the condition that restriction be an equivalence on all coherent diagrams. The localizer axioms express the stability of those constant-coefficient tests under interval homotopy, comma constructions, and integration of fibrations. We complete the reduction from the five displayed axioms to the shorter list, including the role of the twisted-arrow category in proving invariance under opposites. The finite-product and finite-sum questions are then settled by explicit closure arguments and descent of the product and coproduct adjunctions. No model structure is needed for these formal consequences.

Realization is a further theorem. For an accessible fundamental localizer, the later model structures and derivator constructions give a coefficient theory that realizes its weak equivalences \citep[Theorems~4.2.15 and 5.2.15]{Cisinski2006}; \citep[Theorem~6.11 and its dual]{Cisinski2003}. This does not reconstruct an arbitrary original derivator from $W_{\D}$. The represented derivators of sets and of abelian groups have the same component-detecting localizer, although their coefficient categories are inequivalent. The represented truth-value lattice gives an equally concrete nonzero example whose localizer detects only emptiness, disproving any attempt to impose component detection on all nonzero derivators. These examples show precisely what Chapter X has isolated: the homotopy shapes visible to a theory, while its further algebra of coefficients remains additional structure.
\end{otherlanguage}
\dgrchunkend
